\section{Preliminaries for general theory}\label{section2}
\subsection{Root systems}\label{root_sys}
Let $\fg$ be a reductive Lie algebra with Cartan involution $\vartheta$. We decompose $\fg$ into a sum of simple non-compact subalgebras, a semi-simple compact subalgebra and an abelian subalgebra as
\begin{gather*} \fg=\fg_{(1)}\oplus\cdots\oplus\fg_{(m)}\oplus\fg_{\text{cpt}}\oplus\fz(\fg). \end{gather*}
We assume that each simple non-compact subalgebra $\fg_{(i)}$ is of Hermitian type, that is,
its maximal compact subalgebra $\fk_{(i)}:=\fg^\vartheta_{(i)}$ has a 1-dimensional center $\fz(\fk_{(i)})$, and also that the abelian part $\fz(\fg)$ is fixed by $\vartheta$. For each $i$, we fix an element $z_{(i)}\in\fz(\fk_{(i)})$ such that $\operatorname{ad}(z_{(i)})$ has eigenvalues $+\sqrt{-1}$, $0$, $-\sqrt{-1}$,
and decompose the complexified Lie algebra $\fg_{(i)}^\BC$ into eigenspaces under $\operatorname{ad}(z_{(i)})$ as
\begin{gather*} \fg_{(i)}^\BC=\fp^+_{(i)}\oplus \fk^\BC_{(i)}\oplus\fp^-_{(i)}. \end{gather*}
We denote
\begin{alignat*}{3}
& \fp^+ :=\fp^+_{(1)}\oplus\cdots\oplus\fp^+_{(m)},\qquad && \fk^\BC:=\fk_{(1)}^\BC\oplus\cdots\oplus\fk_{(m)}^\BC\oplus\fg_{\text{cpt}}^\BC\oplus\fz(\fg)^\BC,&\\
&\fp^-:=\fp^-_{(1)}\oplus\cdots\oplus\fp^-_{(m)},\qquad && \fk :=\fk_{(1)}\oplus\cdots\oplus\fk_{(m)}\oplus\fg_{\text{cpt}}\oplus\fz(\fg)=\fg^\vartheta,&
\end{alignat*}
so that
\begin{gather*} \fg^\BC=\fp^+\oplus\fk^\BC\oplus\fp^-. \end{gather*}
We denote the anti-holomorphic extension of the Cartan involution $\vartheta$ on $\fg^\BC$ by the same sym\-bol~$\vartheta$.
Also, let $\hat{\vartheta}:=\vartheta\circ \operatorname{Ad}({\rm e}^{\pi z})$ ($z:=\sum_i{z_{(i)}}$) be the anti-holomorphic involution on~$\fg^\BC$ fixing~$\fg$.

Next, we fix a Cartan subalgebra $\fh\subset\fk$. Then $\fh^\BC$ automatically becomes a Cartan subalgebra of $\fg^\BC$.
We set $\fh_{(i)}:=\fh\cap\fg_{(i)}$.
Let $\Delta_{\fg^\BC_{(i)}}=\Delta(\fg^\BC_{(i)},\fh^{\BC}_{(i)})$ be the root system of $\fg^\BC_{(i)}$,
and let $\Delta_{\fp^\pm_{(i)}}$, $\Delta_{\fk^\BC_{(i)}}$ be the set of roots
such that the corresponding root space is contained in $\fp^\pm_{(i)}$, $\fk^\BC_{(i)}$ respectively.
We fix a positive system $\Delta_{\fg^\BC_{(i)},+}\subset \Delta_{\fg^\BC_{(i)}}$
such that $\Delta_{\fp^+_{(i)}}\subset\Delta_{\fg^\BC_{(i)},+}$,
and denote $\Delta_{\fk^\BC_{(i)},+}:=\Delta_{\fk^\BC_{(i)}}\cap \Delta_{\fg^\BC_{(i)},+}$.
Then we can take a system of strongly orthogonal roots
$\{\gamma_{1,(i)},\ldots,\gamma_{r_{(i)},(i)}\}\subset \Delta_{\fp^+_{(i)}}$,
where $r_{(i)}=\rank_\BR\fg_{(i)}$, such that
\begin{enumerate}\itemsep=0pt
\item $\gamma_{1,(i)}$ is the highest root in $\Delta_{\fp^+_{(i)}}$,
\item $\gamma_{k,(i)}$ is the root in $\Delta_{\fp^+_{(i)}}$
which is highest among the roots strongly orthogonal to each $\gamma_{j,(i)}$ with $1\le j\le k-1$.
\end{enumerate}
For each $j$, let $\fp^+_{jj,(i)}$ be the root space corresponding to $\gamma_{j,(i)}$.
We take an element $e_{j,(i)}\in\fp^+_{jj,(i)}$ such that
\begin{gather*} -[[e_{j,(i)},\vartheta e_{j,(i)}],e_{j,(i)}]=2e_{j,(i)}, \end{gather*}
and set
\begin{alignat*}{4}
& h_{j,(i)}:=-[e_{j,(i)},\vartheta e_{j,(i)}]\in\sqrt{-1}\fh_{(i)},\qquad && e_{(i)}:=\sum_{j=1}^{r_{(i)}}e_{j,(i)}\in\fp^+_{(i)},\qquad && e:=\sum_{i=1}^m e_{(i)}\in\fp^+,& \\
&\fa_{\fl,(i)}:=\bigoplus_{j=1}^{r_{(i)}}\BR h_{j,(i)}\subset\sqrt{-1}\fh_{(i)},\qquad && \fa_{(i)}^+:=\bigoplus_{j=1}^{r_{(i)}}\BR e_{j,(i)}\subset\fp^+_{(i)}.&&&
\end{alignat*}
Then the restricted root system $\Sigma=\Sigma\big(\fg^\BC_{(i)},\fa_{\fl,(i)}^\BC\big)$ is one of
\begin{gather*} \Sigma=\big\{ \tfrac{1}{2}(\gamma_{j,(i)}-\gamma_{k,(i)})\big|_{\fa_{\fl,(i)}}\!\!\colon
 1\le j, k\le r_{(i)}, j\ne k \big\}\\
 \hphantom{\Sigma=}{}
\cup\big\{{\pm}\tfrac{1}{2}(\gamma_{j,(i)}+\gamma_{k,(i)})\big|_{\fa_{\fl,(i)}}\!\!\colon 1\le j\le k\le r_{(i)}\big\} \end{gather*}
(type $C_{r_{(i)}}$), or
\begin{gather*} \Sigma=(\text{as above})\cup\big\{ {\pm}\tfrac{1}{2}\gamma_{j,(i)}\big|_{\fa_{\fl,(i)}}\colon 1\le j\le r_{(i)}\big\} \end{gather*}
(type $BC_{r_{(i)}}$). For $1\le j\le k\le r_{(i)}$ we set
\begin{gather*}
\fp^+_{jk,(i)} :=\big\{x\in\fp^+_{(i)}\colon \operatorname{ad}(l)x=\tfrac{1}{2}(\gamma_{j,(i)}+\gamma_{k,(i)})(l)x\text{ for all }l\in\fa_{\fl,(i)}\big\},\\
\fp^+_{0j,(i)} :=\big\{x\in\fp^+_{(i)}\colon \operatorname{ad}(l)x=\tfrac{1}{2}\gamma_{j,(i)}(l)x\text{ for all } l\in\fa_{\fl,(i)}\big\}.
\end{gather*}
Then we have
\begin{gather*} \fp^+_{(i)}=\bigoplus_{\substack{0\le j\le k\le r_{(i)}\\ (j,k)\ne (0,0)}}\fp^+_{jk,(i)}. \end{gather*}
We set
\begin{alignat*}{4}
&\fp^+_{\rT,(i)}:=\bigoplus_{1\le j\le k\le r_{(i)}}\fp^+_{jk,(i)},\qquad &&
\fp^-_{\rT,(i)}:=\vartheta\fp^+_{\rT,(i)},\qquad &&
\fp^+_{\rT}:=\bigoplus_{i=1}^m \fp^+_{\rT,(i)}, &\\
&\fk^\BC_{\rT,(i)}:=[\fp^+_{\rT,(i)},\fp^-_{\rT,(i)}], \qquad && \fk_{\rT,(i)} :=\fk^\BC_{\rT,(i)}\cap \fk_{(i)}, &&&\\
& \fg^\BC_{\rT,(i)}:=\fp^+_{\rT,(i)}\oplus \fk^\BC_{\rT,(i)}\oplus \fp^-_{\rT,(i)},\qquad && \fg_{\rT,(i)}:=\fg^\BC_{\rT,(i)}\cap \fg_{(i)},&&&
\end{alignat*}
and we define the integers
\begin{gather*}
d_{(i)} :=\dim\fp^+_{12,(i)},\qquad b_{(i)}:=\dim\fp^+_{01,(i)},\\
n_{(i)}:=\dim\fp^+_{(i)}=r_{(i)}+\tfrac{1}{2}r_{(i)}(r_{(i)}-1)d_{(i)}+b_{(i)}r_{(i)},\qquad
n:=\dim\fp^+=\sum_{i=1}^m n_{(i)},\\
n_{\rT,(i)}:=\dim\fp^+_{\rT,(i)}=r_{(i)}+\tfrac{1}{2}r_{(i)}(r_{(i)}-1)d_{(i)},\\
p_{(i)}:=2+(r_{(i)}-1)d_{(i)}+b_{(i)}.
\end{gather*}

Throughout the paper, let $G^\BC$ be a connected complex Lie group with Lie algebra $\fg^\BC$, and let
$G$, $K^\BC$, $K$, $G^\BC_{(i)}$, $G_{(i)}$, $K^\BC_{(i)}$, $K_{(i)}$, $G^\BC_{\rT,(i)}$, $G_{\rT,(i)}$, $K^\BC_{\rT,(i)}$, $K_{\rT,(i)}$
be the connected Lie subgroup with Lie algebras
$\fg$, $\fk^\BC$, $\fk$, $\fg^\BC_{(i)}$, $\fg_{(i)}$, $\fk^\BC_{(i)}$, $\fk_{(i)}$,
$\fg^\BC_{\rT,(i)}$, $\fg_{\rT,(i)}$, $\fk^\BC_{\rT,(i)}$, $\fk_{\rT,(i)}$ respectively.
Also, let
\begin{gather*} K_{L,(i)}:=\big\{ k\in K_{\rT,(i)}\colon \operatorname{Ad}(k)e_{(i)}=e_{(i)}\big\}, \end{gather*}
which is possibly non-connected, and we denote its Lie algebra by $\fk_{\fl,(i)}$.

For $k\in K^\BC$, we write $k^*:=(\vartheta k)^{-1}$.
Then for each $i$, there exists a unique Hermitian inner product $(\cdot|\cdot)_{\fp^+_{(i)}}$,
holomorphic in the first variable and anti-holomorphic in the second variable, such that
\begin{gather*}
(\operatorname{Ad}(k)x|y)_{\fp^+_{(i)}} =(x|\operatorname{Ad}(k^*)y)_{\fp^+_{(i)}}, \qquad x,y\in\fp^+_{(i)},\quad k\in K^\BC_{(i)},\qquad
(e_{1,(i)}|e_{1,(i)})_{\fp^+_{(i)}} =1.
\end{gather*}
This is proportional to the restriction of the Killing form of $\fg_{(i)}^\BC$ on $\fp^+_{(i)}\times \fp^-_{(i)}$,
if we identify $\fp^+_{(i)}$ and $\fp^-_{(i)}$ through $\vartheta$. By summing these inner products, we define
\begin{gather}\label{innerprod}
(x|y)=(x|y)_{\fp^+}:=\sum_{i=1}^m (x_i|y_i)_{\fp^+_{(i)}} \qquad
 x=\sum_{i=1}^m x_i, \qquad y=\sum_{i=1}^m y_i\in\fp^+=\bigoplus_{i=1}^m\fp^+_{(i)} .
\end{gather}
From now on we omit $\operatorname{Ad}$ or $\operatorname{ad}$ if there is no confusion, so that $(kx|y)_{\fp^+}=(x|k^*y)_{\fp^+}$.

\subsection{Operations on Jordan triple systems}
$\fp^+$ has a Hermitian positive Jordan triple system structure with the product
\begin{gather*} (x,y,z)\mapsto -[[x,\vartheta y],z]. \end{gather*}
For $x,y\in\fp^+$, let $D(x,y)$ be the linear map, $Q(x,y)$ be the anti-linear map on $\fp^+$ defined by
\begin{gather*} D(x,y):=-\operatorname{ad}([x,\vartheta y])\bigr|_{\fp^+},\qquad Q(x,y):=\operatorname{ad}(x)\operatorname{ad}(y)\vartheta\bigr|_{\fp^+}, \end{gather*}
and let $Q(x):=\frac{1}{2}Q(x,x)$.
We recall that, for $x,y\in\fp^+$, the \textit{Bergman operator} $B(x,y)\in\End(\fp^+)$ is defined as
\begin{gather*} B(x,y):=I-D(x,y)+Q(x)Q(y)\in\End(\fp^+). \end{gather*}
We say $(x,y)\in\fp^+\times\fp^+$ is \textit{quasi-invertible} if $B(x,y)$ (or equivalently $B(y,x)$) is invertible,
and in this case the \textit{quasi-inverse} $x^y$ is defined as
\begin{gather*} x^y:=B(x,y)^{-1} \left(x-Q(x)y\right)\in\fp^+. \end{gather*}
Then if $B(x,y)$ is invertible, then there exists an element $k\in K^\BC$ such that $B(x,y)z=\operatorname{Ad}(k)z$ holds for any $z\in\fp^+$.
Also, $B(x,y)$ and $x^y$ satisfy the following properties. For $x,y,z\in\fp^+$ and $k\in K^\BC$, if $(x,y)$ is quasi-invertible, then
\begin{alignat}{3}
& B(kx,k^{*-1}y) =kB(x,y)k^{-1}, &&\label{Bergman_k}\\
& B(x,y)B(x^y,z) =B(x,y+z) \qquad & & \text{\cite[Part V, Proposition III.3.1, (J6.4)]{FKKLR}},&\label{Bergman_left}\\
& B(z,x^y)B(y,x) =B(y+z,x) \qquad & & \text{\cite[Part V, Proposition III.3.1, (J6.4$'$)]{FKKLR}},&\label{Bergman_right}\\
& (kx)^{k^{*-1}y} =k(x^y), &&& \label{quasiinv_k}\\
& x^{y+z} =(x^y)^z \qquad & & \text{\cite[Part V, Theorem III.5.1(i)]{FKKLR}},& \label{quasiinv_add}\\
&(x+z)^y=x^y+B(x,y)^{-1}z^{(y^x)} \qquad & & \text{\cite[Part V, Theorem III.5.1(ii)]{FKKLR}}& \label{quasiinv_twice}
\end{alignat}
hold. Here, the equality (\ref{quasiinv_add}) holds when one of $(x,y+z)$ or $(x^y,z)$ is quasi-invertible,
and the other also becomes quasi-invertible. Similarly, the equality~(\ref{quasiinv_twice}) holds when one of $(x+z,y)$ or $(z,y^x)$ is
quasi-invertible, and then the other also is. Also, for the Bergman operator, we can show directly from the definition that,
if $\fp^+_1,\fp^+_2\subset \fp^+$ are Jordan triple subsystems such that $D(\fp^+_1,\fp^+_2)=\{0\}$ (we do not assume they are ideals),
then we have
\begin{gather}\label{Bergman_decomp}
B(x_1+x_2,y_1+y_2)=B(x_1,y_1)B(x_2,y_2), \qquad x_1,y_1\in\fp^+_1, \quad x_2,y_2\in\fp^+_2.
\end{gather}

Next, we recall the spectral decomposition and the spectral norm. For any $x_i\in\fp^+_{(i)}$, there exist complex numbers
$a_{1,i},\ldots, a_{r_{(i)},i}$ and an element $k_i\in K_{(i)}$ such that
\begin{gather*} x_i=k_i\sum_{j=1}^{r_{(i)}} a_{j,i}e_{j,(i)}. \end{gather*}
The set $\{a_{1,i},\ldots,a_{r_{(i)},i}\}$ is unique under the condition that $a_{j,i}\in\BR_{\ge 0}$ and
$a_{1,i}\ge\cdots\ge a_{r_{(i)},i}\ge 0$. This is called the \textit{spectral decomposition}.
For $x=\sum\limits_{i=1}^m x_i=\sum\limits_{i=1}^m k_i\sum\limits_{j=1}^{r_{(i)}}a_{j,i}e_{j,(i)}\in\fp^+=\bigoplus\limits_{i=1}^m \fp^+_{(i)}$,
the \textit{spectral norm} is defined as
\begin{gather}\label{spectral_norm}
|x|_\infty=|x|_{\fp^+,\infty}:=\max_{1\le i\le m}\max_{1\le j\le r_{(i)}}|a_{j,i}|.
\end{gather}
In fact this becomes a norm on the vector space $\fp^+$ (see \cite[Part V, Proposition VI.4.1]{FKKLR}).

Next, for each $i$, let $h_{(i)}(x,y)\in\cP(\fp^+\times\overline{\fp^+})$ be the \textit{generic norm} on $\fp^+_{(i)}$.
This is the polynomial, holomorphic in $x$ and anti-holomorphic in $y$, satisfying
\begin{gather*} \Det_{\fp^+_{(i)}}(B(x_i,y_i))=h_{(i)}(x_i,y_i)^{p_{(i)}}, \qquad x_i,y_i\in\fp^+_{(i)} . \end{gather*}
If $x_i=\sum\limits_{j=1}^{r_{(i)}}a_je_{j,(i)}$, $y_i=\sum\limits_{j=1}^{r_{(i)}}b_je_{j,(i)}\in\fa^+_{(i)}\subset\fp^+_{(i)}$,
then $h_{(i)}(x_i,y_i)$ is given by
\begin{gather*} h_{(i)}(x_i,y_i)=\prod_{j=1}^{r_{(i)}}(1-a_j\overline{b_j}). \end{gather*}
For later use we abbreviate
\begin{gather*} \Det_{\fp^+}(B(x,y))^{-1}=\prod_{i=1}^m h_{(i)}(x_i,y_i)^{-p_{(i)}}=:h(x,y)^{-p}. \end{gather*}
Also, we abbreviate $B(x,x)=:B(x)$, $h_{(i)}(x_i,x_i)=h_{(i)}(x_i)$. Let
\begin{align}
D\colon \hspace{-3pt}&=(\text{connected component of }\{x\in \fp^+\colon B(x)\text{ is positive definite.}\}\text{ which contains }0) \notag \\
&=\big\{x\in \fp^+\colon |x|_\infty <1\}=\{x\in \fp^+\colon |D(x,x)|_{\fp^+,\mathrm{op}}<2\big\}\label{BSD}
\end{align}
be the \textit{bounded symmetric domain}, which is diffeomorphic to $G/K$ via the Borel embedding which we will review later
(for these equalities see \cite[Part~V, Proposition~VI.4.2]{FKKLR}. Here $|\cdot|_{\fp^+,\mathrm{op}}$ denotes the
operator norm on $\End(\fp^+)$ with respect to $|\cdot|_{\fp^+}$).
Then if $x,y\in D$, $B(x,y)$ is invertible, and thus it is in the image of $K^\BC$.
Moreover, since $D$ is simply connected, there exists a holomorphic map $\tilde{B}\colon D\times\overline{D}\to K^\BC$
(or $\tilde{B}\colon D\times\overline{D}\to \tilde{K}^\BC$, where $\tilde{K}^\BC$ is the universal covering group of $K^\BC$)
such that
\begin{gather*} \operatorname{Ad}(\tilde{B}(x,y))=B(x,y)\in\End(\fp^+), \qquad \tilde{B}(0,0)=\bone_{K^\BC}\in K^\BC \quad \big(\text{resp.} \ {} \in \tilde{K}^\BC\big) \end{gather*}
holds. From now on we omit the tilde, and use the same symbol $B$ instead of $\tilde{B}$.

Next we consider $\fp^+_\rT$. This has a complex Jordan algebra structure with the product
\begin{gather*} (x,y)\mapsto x\cdot y:=-\tfrac{1}{2}[[x,\vartheta e],y]. \end{gather*}
We recall the quadratic map $P\colon \fp^+_\rT\to\End(\fp^+_\rT)$ defined by
\begin{gather}\label{quad_repn}
P(x)y:=2x\cdot(y\cdot x)-y\cdot(x\cdot x)=Q(x)Q(e)y, \qquad x,y\in\fp^+_\rT .
\end{gather}
If $y$ is in the real form $\left\{y\in\fp^+_\rT\colon Q(e)y=y\right\}$ of $\fp^+_\rT$,
then $P(x)y=-\frac{1}{2}[[x,\vartheta y],x]=Q(x)y$ holds.
Next we review the \textit{determinant polynomials} on Jordan algebras.
On each simple component~$\fp^+_{\rT,(i)}$ there exists a determinant polynomial $\Delta_{(i)}$,
which is the homogeneous polynomial of degree~$r_{(i)}$ satisfying
\begin{gather*}
\Delta_{(i)}(kx)=\Delta_{(i)}(ke_{(i)})\Delta_{(i)}(x)
\qquad \text{for all} \ k\in K^\BC_{\rT,(i)},\ x\in \fp^+_{\rT,(i)},\qquad
\Delta_{(i)}(e_{(i)})=1.
\end{gather*}
The quadratic map $P$ and the determinant polynomials are related as
\begin{gather*} \Det_{\fp^+_{\rT,(i)}}(P(x_i))=\Delta(x_i)^{2n_{\rT,(i)}/r_{(i)}}, \qquad x_i\in\fp^+_{\rT,(i)} . \end{gather*}
We \looseness=-1 extend $\Delta_{(i)}$ on $\fp^+_{(i)}$ such that it does not depend on $\big(\fp^+_{\rT,(i)}\big)^\bot=\bigoplus\limits_{j=1}^{r_{(i)}} \fp^+_{0j,(i)}$, and denote by the same symbol $\Delta_{(i)}$.
Then the determinant polynomial $\Delta_{(i)}$ and the generic norm $h_{(i)}$ are related as
\begin{gather}\label{det_norm}
\Delta_{(i)}(e_{(i)}-x)=h_{(i)}(x,e_{(i)}), \qquad x\in\fp^+_{(i)}.
\end{gather}
For the theory of Jordan algebras and Jordan triple systems, see, e.g., \cite[Part V]{FKKLR}, \cite{FK, L,Sat}.

\subsection{Polynomials on Jordan triple systems}
Let $\cP(\fp^+)$ be the space of all holomorphic polynomials on $\fp^+$. Then $K^\BC$ acts on $\cP(\fp^+)$ by
\begin{gather*} (\operatorname{Ad}|_{\fp^+})^\vee(k)f(x):=f\big(k^{-1}x\big), \qquad k\in K^\BC,\quad f\in\cP(\fp^+). \end{gather*}
Then clearly we have $\cP(\fp^+)\simeq \cP(\fp^+_{(1)})\otimes\cdots\otimes\cP(\fp^+_{(m)})$, according to the
simple decomposition of the Jordan triple system $\fp^+=\fp^+_{(1)}\oplus\cdots\oplus\fp^+_{(m)}$.
In the rest of this subsection, we assume $\fg$ is simple, and we drop the subscript $(i)$. We set
\begin{gather*} \BZ_{++}^r:=\big\{\bm=(m_1,\ldots,m_r)\in\BZ^r\colon m_1\ge \cdots \ge m_r\ge 0\big\}. \end{gather*}
Then $\cP(\fp^+)$ is decomposed as follows.
\begin{Theorem}[Hua--Kostant--Schmid, {\cite[Part~III, Theorem~V.2.1]{FKKLR}}]\label{HKS}
Under the $K^\BC$-action, $\cP(\fp^+)$ is decomposed as
\begin{gather*} \cP(\fp^+)=\bigoplus_{\bm\in\BZ^{r}_{++}}\cP_\bm(\fp^+), \end{gather*}
where $\cP_\bm(\fp^+)$ is the irreducible representation of $K^\BC$ with lowest weight
$-m_1\gamma_1-\cdots-m_r\gamma_r$. Moreover, each $\cP_\bm(\fp^+)$ has a nonzero $K_L$-invariant polynomial,
which is unique up to scalar multiple.
\end{Theorem}
Let $d_\bm^{(d,r,b)}:=\dim\cP_\bm(\fp^+)$, $d_\bm^{(d,r,0)}:=\dim\cP_\bm(\fp^+_\rT)$,
where $\cP_\bm(\fp^+_\rT)$ is the irreducible representation of $K^\BC_\rT$ with lowest weight $-m_1\gamma_1-\cdots-m_r\gamma_r$,
and let $\Phi_\bm^{(d,r)}$ be the $K_L$-invariant polynomial in $\cP_\bm(\fp^+)$
such that $\Phi_\bm^{(d,r)}(e)=1$. Especially, when $\bm=(m,\ldots,m)$, then $\Phi_{(m,\ldots,m)}^{(d,r)}(x)=\Delta(x)^m$ holds.

Next we recall the \textit{Fischer inner product}.
For two holomorphic polynomials $f,g\in\cP(\fp^+)$, it is defined as
\begin{gather}\label{Fischer}
\langle f,g\rangle_F:=\frac{1}{\pi^n}\int_{\fp^+}f(x)\overline{g(x)} {\rm e}^{-|x|_{\fp^+}^2}{\rm d}x.
\end{gather}
This integral converges for polynomials $f$ and $g$, and the reproducing kernel is given by ${\rm e}^{(x|y)_{\fp^+}}$.
Let $\bK_\bm^{(d)}(x,y)\in\cP(\fp^+\times\overline{\fp^+})$ be the reproducing kernel of $\cP_\bm(\fp^+)$
with respect to $\langle\cdot,\cdot\rangle_F$, so that $\sum\limits_{\bm\in\BZ_{++}^r}\bK_\bm^{(d)}(x,y)={\rm e}^{(x|y)_{\fp^+}}$. Then the following holds.
\begin{Proposition}[{\cite[Part III, Lemma V.3.1(a), Theorem V.3.4]{FKKLR}}]\label{Kmxe}
\begin{gather*} \bK_\bm^{(d)}(x,e)=\frac{d_\bm^{(d,r,b)}}{\left(\frac{n}{r}\right)_{\bm,d}}\Phi_\bm^{(d,r)}(x)
=\frac{d_\bm^{(d,r,0)}}{\left(\frac{n_\rT}{r}\right)_{\bm,d}}\Phi_\bm^{(d,r)}(x). \end{gather*}
\end{Proposition}
Here, for $\lambda\in\BC$, $\bs\in\BC^r$, $\bm\in(\BZ_{\ge 0})^r$ and $d\in\BZ_{\ge 0}$, $(\lambda+\bs)_{\bm,d}$ is defined as
\begin{gather}\label{Pochhammer}
(\lambda+\bs)_{\bm,d}:=\prod_{j=1}^r\left(\lambda+s_j-\frac{d}{2}(j-1)\right)_{m_j},\qquad
(\lambda)_m:=\lambda(\lambda+1)\cdots(\lambda+m-1),
\end{gather}
and we write $(\lambda+(0,\ldots,0))_{\bm,d}=(\lambda)_{\bm,d}$.
We renormalize $\Phi_\bm^{(d,r)}$ as
\begin{gather}\label{Phitilde1}
\tilde{\Phi}_\bm^{(d)}(x):=\frac{d_\bm^{(d,r,b)}}{\left(\frac{n}{r}\right)_{\bm,d}}\Phi_\bm^{(d,r)}(x)
=\frac{d_\bm^{(d,r,0)}}{\left(\frac{n_\rT}{r}\right)_{\bm,d}}\Phi_\bm^{(d,r)}(x),
\end{gather}
so that
\begin{gather*} {\rm e}^{(x|e)_{\fp^+}}=\sum_{\bm\in\BZ_{++}^r}\bK_\bm^{(d)}(x,e)=\sum_{\bm\in\BZ_{++}^r}\tilde{\Phi}_\bm^{(d)}(x). \end{gather*}
For example, when $\fp^+=M(r,\BC)$ (i.e., $G=\operatorname{SU}(r,r)$), if the eigenvalues of $x\in M(r,\BC)$ are $t_1,\ldots,t_r$, we have
\begin{gather}
\tilde{\Phi}_\bm^{(2)}(x) = \!\left(\!\frac{\prod\limits_{i<j}(m_i-m_j-i+j)}{\prod\limits_{i=1}^r (r-i)!}\!\right)^2\!\!\! \frac{1}{\prod\limits_{i=1}^r (r-i+1)_{m_i}}
 \frac{\prod\limits_{i=1}^r (r-i)!}{\prod\limits_{i<j}(m_i-m_j-i+j)}
\frac{\det\big(\big(t_i^{m_j+r-j}\big)_{i,j}\big)}{\det\big(\big(t_i^{r-j}\big)_{i,j}\big)}\notag \\
\hphantom{\tilde{\Phi}_\bm^{(2)}(x)}{} =\frac{\prod\limits_{i<j}(m_i-m_j-i+j)}{\prod\limits_{i=1}^r (m_i+r-i)!}
\frac{\det\left((t_i^{m_j+r-j})_{i,j}\right)}{\det\left((t_i^{r-j})_{i,j}\right)}.\label{Schur}
\end{gather}
Then $\tilde{\Phi}_\bm^{(d)}(x)$ does not depend on $r$ in the following sense.
Since $\tilde{\Phi}_\bm^{(d)}$ is $K_L$-invariant, it is determined by the value on $\fa^+\subset\fp^+$.
Thus for $x=a_1e_1+\cdots+a_re_r\in\fa^+$, we write
\begin{gather*} \tilde{\Phi}_\bm^{(d)}(x)=:\tilde{\Phi}_\bm^{(d)}(a_1,\ldots,a_r). \end{gather*}
Then this does not depend on $r$, that is,
\begin{gather*} \tilde{\Phi}_\bm^{(d,r)}(a_1,\ldots,a_{r-1},0)=\tilde{\Phi}_\bm^{(d,r-1)}(a_1,\ldots,a_{r-1}) \end{gather*}
holds. Also, when $\fp^+=\Sym(r,\BC)$, $M(q,s;\BC)$ or $\Skew(s,\BC)$ (i.e., $G=\operatorname{Sp}(r,\BR)$, $\operatorname{SU}(q,s)$ or $\operatorname{SO}^*(2s)$ respectively),
for $x,y\in\fp^+$, $\bK_\bm^{(d)}(x,y)$ depends only on the eigenvalues of $xy^*$, so following the notation in \cite{M} we write
\begin{gather}\label{Phitilde2}
\bK_\bm^{(d)}(x,y)=:\tilde{\Phi}_\bm^{(d)}(xy^*)=\tilde{\Phi}_\bm^{(d)}(y^*x).
\end{gather}


\subsection{Holomorphic discrete series representations}\label{HDS}
In this subsection we recall the explicit realization of the holomorphic discrete series representation
of the universal covering group $\tilde{G}$. First we recall the Borel embedding,
\begin{gather*} \xymatrix{ G/K \ar[r] \ar@{-->}[d]^{\mbox{\rotatebox{90}{$\sim$}}} & G^\BC/K^\BC P^- \\
 D \ar@{^{(}->}[r] & \fp^+, \ar[u]_{\exp} } \end{gather*}
where $P^\pm:=\exp(\fp^\pm)$. When $g\in G^\BC$ and $x\in \fp^+$ satisfy $g\exp(x)\in P^+K^\BC P^-$, we write
\begin{gather}\label{cocycle}
g\exp(x)=\exp(\pi^+(g,x))\kappa(g,x)\exp(\pi^-(g,x)),
\end{gather}
where $\pi^+(g,x)\in\fp^+$, $\kappa(g,x)\in K^\BC$, and $\pi^-(g,x)\in\fp^-$.
If $g=k\in K^\BC$, $g=\exp(y)\in P^+$ or $g=\exp(\vartheta y)\in P^-$ with $y\in\fp^+$, we have
\begin{alignat*}{3}
&\pi^+(k,x)=kx,\qquad && \kappa(k,x)=k,&\\
& \pi^+(\exp(y),x)=x+y,\qquad && \kappa(\exp(y),x) =\bone_{K^\BC},& \\
& \pi^+(\exp(\vartheta y),x) =x^y,\qquad && \operatorname{Ad}(\kappa(\exp(\vartheta y),x))|_{\fp^+}=B(x,y)^{-1}.&
\end{alignat*}
$\pi^+$ gives the birational action of $G^\BC$ on $\fp^+$, and from now on we abbreviate $\pi^+(g,x)=:gx$.
Especially, if $x\in D$ and $g\in G$, then automatically $gx\in D$ and $\kappa(g,x)$ is well-defined,
and the action of $G$ on $D$ is transitive. Since $D$ is simply connected, the map $\kappa\colon G\times D\to K^\BC$
lifts to the universal covering space, that is, $\kappa\colon \tilde{G}\times D\to \tilde{K}^\BC$ is well-defined.
We denote this extended map by the same symbol~$\kappa$. Then for $x,y\in\fp^+$ and $g\in G^\BC$,
\begin{gather}\label{Bergman_property}
B\big(gx,\big(\hat{\vartheta}g\big)y\big)=\kappa(g,x)B(x,y)\kappa\big(\hat{\vartheta}g,y\big)^*
\end{gather}
holds in $\End(\fp^+)$, where $\hat{\vartheta}$ is the anti-holomorphic involution of $G^\BC$ fixing $G$, and $\operatorname{Ad}$ is omitted.
If $g\in G$ (i.e., $g=\hat{\vartheta}g$) and $x,y\in D$, this also holds in $K^\BC$, regarding $B(x,y)$ as the element of~$K^\BC$.
This formula is also verified in $\tilde{K}^\BC$ if $g\in \tilde{G}$.

Now let $(\tau,V)$ be an irreducible holomorphic representation of $\tilde{K}^\BC$ with $\tilde{K}$-invariant inner product $(\cdot,\cdot)_\tau$.
We consider the space of holomorphic sections of the homogeneous vector bundle on $G/K$ with fiber~$V$. Then since $D\simeq G/K$ is contractible, it is isomorphic to the space of $V$-valued holomorphic functions on~$D$
\begin{gather*} \Gamma_\cO\big(G/K, \tilde{G}\times_{\tilde{K}}V\big)\simeq \cO(D,V). \end{gather*}
Via this identification, $\tilde{G}$ acts on $\cO(D,V)=\cO_\tau(D,V)$ by
\begin{gather*} \hat{\tau}(g)f(x)=\tau\big(\kappa\big(g^{-1},x\big)\big)^{-1}f\big(g^{-1}x\big),\qquad g\in\tilde{G}, \quad x\in D, \quad f\in\cO(D,V) , \end{gather*}
and the function $\tau(B(x,y))\in \cO(D\times\overline{D},\End(V))$ is invariant under the diagonal action of~$\tilde{G}$.
If the function $\tau(B(x,y))$ is positive-definite, that is,
$\sum\limits_{j,k=1}^N(\tau(B(x_j,x_k))v_j,v_k)_\tau\ge 0$ holds for any $\{x_j\}_{j=1}^N\subset D$ and $\{v_j\}_{j=1}^N\subset V$,
then there exists a unique Hilbert subspace $\cH_\tau(D,V)\subset\cO_\tau(D,V)$ with the reproducing kernel $\tau(B(x,y))$,
on which $\tilde{G}$ acts unitarily via $\hat{\tau}$.
This representation $(\hat{\tau},\cH_\tau(D,V))$ is called a unitary highest weight representation.
Especially, if its inner product is given by the converging integral
\begin{gather*} \langle f,g\rangle_{\hat{\tau}}:=C_\tau\int_D \big(\tau\big(B(x)^{-1}\big)f(x),g(x)\big)_{\tau}h(x)^{-p}{\rm d}x, \end{gather*}
where $h(x)^{-p}{\rm d}x:=\prod\limits_{i=1}^m h_{(i)}(x_i)^{-p_{(i)}}{\rm d}x=\Det(B(x))^{-1}{\rm d}x$ is the $G$-invariant measure on~$D$,
and~$C_\tau$ is the constant such that $\Vert v\Vert_{\hat{\tau}}=|v|_\tau$ holds for any constant functions
(or elements in the minimal $K$-type)~$v$,
then $(\hat{\tau},\cH_\tau(D,V))$ is called a holomorphic discrete series representation.
In this case, all bounded holomorphic functions on $D$ belong to $\cH_\tau(D,V)$,
and especially the space of $\tilde{K}$-finite vectors is equal to the space of all polynomials,
\begin{gather*} \cH_\tau(D,V)_{\tilde{K}}=\cO_\tau(D,V)_{\tilde{K}}=\cP(\fp^+,V). \end{gather*}
For general $\cH_\tau(D,V)$ such that the above integral does not converge for any non-zero function,
it may happen that the $\tilde{K}$-finite part $\cH_\tau(D,V)_{\tilde{K}}$ is strictly smaller than $\cP(\fp^+,V)$.

Now we assume $G$ is simple. Let $\chi$ be the character of $\tilde{K}^\BC$ such that
\begin{gather}\label{char}
\chi(k)^p=\Det(\operatorname{Ad}(k)|_{\fp^+}),\qquad \text{or} \qquad \chi(B(x,y))=h(x,y).
\end{gather}
Then for $x,y\in\fp^+$ we have ${\rm d}\chi([x,-\vartheta y])=(x|y)_{\fp^+}$, where $(\cdot|\cdot)_{\fp^+}$ is as in (\ref{innerprod}).
Let $(\tau_0,V)$ be a fixed irreducleble representation of $K^\BC$.
Then for $\lambda\in\BR$, $(\tau,V)=(\tau_0\otimes\chi^{-\lambda},V)$ is again a~representation of $\tilde{K}^\BC$.
In this case we denote $\cH_\tau(D,V)=:\cH_\lambda(D,V)$.
If $\lambda$ is sufficiently large, this becomes a holomorphic discrete series representation.
The parameter $\lambda$ such that the unitary subrepresentation $\cH_\lambda(D,V)\subset\cO_\lambda(D,V)$ exists is classified by
Enright--Howe--Wallach~\cite{EHW} and Jakobsen~\cite{J}.

When $\cH_\lambda(D,V)$ is holomorphic discrete, we consider the irreducible decomposition of \linebreak
$\cH_\lambda(D,V)_{\tilde{K}}\simeq \cP(\fp^+,V)\otimes \chi^{-\lambda}$ as $\tilde{K}^\BC$-modules,
\begin{gather*} \cP(\fp^+,V)\otimes \chi^{-\lambda}\simeq \bigoplus_m W_m\otimes\chi^{-\lambda}, \end{gather*}
such that its components are orthogonal to one another with respect to the Fischer inner product
\begin{gather*} \langle f,g\rangle_{F,\tau_0}:=\frac{1}{\pi^n}\int_{\fp^+} (f(x),g(x) )_{\tau_0}{\rm e}^{-|x|_{\fp^+}^2}{\rm d}x. \end{gather*}
Then since both $\langle \cdot,\cdot\rangle_{F,\tau_0}$ and
$\langle\cdot,\cdot\rangle_{\hat{\tau}}=\langle \cdot,\cdot\rangle_{\lambda,\tau_0}$ are $\tilde{K}$-invariant,
there exists a constant $p_m(\lambda)>0$ such that $\Vert f_m\Vert^2_{\lambda,\tau_0}/\Vert f_m\Vert^2_{F,\tau_0}=p_m(\lambda)$
for any $f_m\in W_m$.
Now we additionally assume that
\begin{gather}\label{assumption_orth}
W_m\perp W_n \text{ in } \langle\cdot,\cdot\rangle_{F,\tau_0}\qquad \text{implies}\qquad
W_m\perp W_n \text{ in } \langle\cdot,\cdot\rangle_{\lambda,\tau_0}
\end{gather}
for sufficiently large $\lambda$.
This holds if $\cP(\fp^+,V)$ is multiplicity-free, or $G=U(q,s)$ and one of~$U(q)$ and~$U(s)$ acts trivially on $V$.
Then $\Vert \cdot\Vert_{\lambda,\tau_0}$ is computed as
\begin{gather}\label{norm_compare}
\Vert f\Vert_{\lambda,\tau_0}^2=\sum_{m,n}\langle f_m,f_n\rangle_{\lambda,\tau_0}=\sum_m p_m(\lambda)\Vert f_m\Vert_{F,\tau_0}^2,
\end{gather}
where $f_m$ is the orthogonal projection of $f$ onto $W_m$, that is, the cross terms vanish.
Especially, if the reproducing kernel with respect to $\langle\cdot,\cdot\rangle_{F,\tau_0}$ is expanded as
\begin{gather*} {\rm e}^{(x|y)_{\fp^+}}I_V=\sum_m \bK_m(x,y)\in\cO\big(\fp^+\times\overline{\fp^+},\End(V)\big), \end{gather*}
where \looseness=-1 $\bK_m(x,y)\in W_m\otimes\overline{W_m}$, then the reproducing kernel with respect to $\langle\cdot,\cdot\rangle_{\lambda,\tau_0}$ is expanded as
\begin{gather}\label{kernel_expansion}
h(x,y)^{-\lambda}\tau_0(B(x,y))=\sum_m p_m(\lambda)^{-1}\bK_m(x,y)\in\cO\big(D\times\overline{D},\End(V)\big).
\end{gather}
Each $p_m(\lambda)$ is meromorphically continued for all $\lambda\in\BC$, and defines a positive definite kernel function if
$p_m(\lambda)^{-1}\ge 0$ for all $m$, and for such $\lambda$ there exists a unitary subrepresentation $\cH_\lambda(D,V)\subset\cO(D,V)$.
Also, if $f_m\in W_m$ we have
\begin{gather}\label{exp_onD}
\langle f_m,{\rm e}^{(\cdot|y)_{\fp^+}}\rangle_{\lambda,\tau_0}=p_m(\lambda)\langle f_m,\bK_m(\cdot,y)\rangle_{F,\tau}=p_m(\lambda)f_m(y).
\end{gather}
Under this assumption (\ref{assumption_orth}), $p_m(\lambda)$ are explicitly computed in \cite{N2} for classical $G$.
We note that if we drop the assumption (\ref{assumption_orth}), then the formulas (\ref{norm_compare}) and (\ref{kernel_expansion}) become
more complicated, and their explicit formulas are not known so far.
Especially if $(\tau_0,V)$ is trivial,
then by \cite[Part~III, Corollary~V.3.9]{FKKLR}, \cite{FK0} for $f_\bm\in\cP_\bm(\fp^+)$ we have
\begin{gather}\label{scalar_norm}
p_\bm(\lambda)=\frac{\Vert f_\bm\Vert_\lambda^2}{\Vert f_\bm\Vert_F^2}=\frac{1}{(\lambda)_{\bm,d}},
\end{gather}
where $(\lambda)_{\bm,d}$ is as (\ref{Pochhammer}), and therefore we have
\begin{gather}\label{ext_of_h}
h(x,y)^{-\lambda}=\sum_{\bm\in\BZ_{++}^r}(\lambda)_{\bm,d}\bK_\bm^{(d)}(x,y),
\end{gather}
where $\bK_\bm^{(d)}(x,y)\in\cP_\bm(\fp^+)\otimes \overline{\cP_\bm(\fp^+)}$ is as in the previous subsection.

\section[Intertwining operators between holomorphic discrete series representations]{Intertwining operators between holomorphic discrete series\\ representations}\label{section3}
\subsection{Setting}
Let $G$ be a connected real reductive Lie group such that each simple non-compact component is of Hermitian type, as in Section~\ref{root_sys},
and let $z\in\fz(\fk)$ be the element such that $\operatorname{ad}(z)|_{\fp^+}=\sqrt{-1}I_{\fp^+}$.
Let $G_1\subset G$ be a connected reductive subgroup which is stable under the Cartan involution $\vartheta$ of $G$.
We denote the Lie algebra of $G_1$ and its Cartan decomposition under $\vartheta$ by $\fg_1=\fk_1\oplus \fp_1$. We assume
\begin{gather}\label{assumption}
z\in \fg_1.
\end{gather}
We set $\fp_1^+:=\fg_1^\BC\cap \fp^+$, $\fp_1^-:=\fg_1^\BC\cap \fp^-$, so that
$\fg_1^\BC=\fp^+_1\oplus \fk^\BC_1\oplus \fp^-_1$.
Also, let $\fp_2^+\subset\fp^+$ be the orthogonal complement of $\fp^+_1$ with respect to the inner product $(\cdot|\cdot)_{\fp^+}$
defined in (\ref{innerprod}).
We define another inner product $(\cdot|\cdot)_{\fp^+_1}$ on $\fp^+_1$ as in (\ref{innerprod}), changing $\fg$ to $\fg_1$,
and let $D_1\subset\fp^+_1$ be the bounded symmetric domain, defined as in (\ref{BSD}). Then we have
\begin{Proposition}\label{BSD_intersection}
$D_1=D\cap \fp^+_1$.
\end{Proposition}
\begin{proof}$D_1=G_1.0\subset D\cap\fp^+_1$ is clear. For $x,y\in \fp^+_1$, let $B_1(x,y)$ be the Bergman operator on~$\fp^+_1$.
Then $B_1(x,y)=B(x,y)|_{\fp^+_1}$ holds. Therefore if $x\in D\cap \fp^+_1$, then $B(x)$ is positive definite,
and hence $B_1(x)=B(x)|_{\fp^+_1}$ is also positive definite. Since $D\cap \fp^+_1$ is connected, $D\cap\fp^+_1\subset D_1$ holds.
\end{proof}

This implies that the spectral norms (\ref{spectral_norm}) on $\fp^+$ and $\fp^+_1$ coincide.
\begin{Corollary}
For $x\in\fp^+_1\subset\fp^+$, $|x|_{\fp^+,\infty}=|x|_{\fp^+_1,\infty}$ holds.
\end{Corollary}
\begin{proof}This follows from $|x|_{\fp^+,\infty}=\inf\left\{t>0\colon \frac{1}{t}x\in D\right\}$,
$|x|_{\fp^+_1,\infty}=\inf\left\{t>0\colon \frac{1}{t}x\in D_1\right\}$ \linebreak (see~(\ref{BSD})) and Proposition~\ref{BSD_intersection}.
\end{proof}

Let $(\tau,V)$ be a representation of $\tilde{K}^\BC$, and consider the representation $(\hat{\tau},\cH_\tau(D,V))$ of $\tilde{G}$, as in Section~\ref{HDS}.
We want to discuss the restriction $\cH_\tau(D,V)|_{\tilde{G}_1}$.
Then since it is discretely decomposable, the space of $\tilde{K}_1$-finite vectors coincides with the space of $\tilde{K}$-finite vectors
(see \cite[Theorem~4.5]{K1}),
which is a subspace of the space of $V$-valued polynomials on $\fp^+$.
\begin{gather*} \cH_\tau(D,V)_{\tilde{K}_1}=\cH_\tau(D,V)_{\tilde{K}}\subset\cP(\fp^+,V). \end{gather*}
Since $\fp^+$ acts on $\cH_\tau(D,V)_{\tilde{K}}\subset\cP(\fp^+,V)$ by 1st order differential operators with constant coefficients,
every $(\fg_1,\tilde{K}_1)$-submodule in $\cH_\tau(D,V)_{\tilde{K}_1}$ has $\fp^+_1$-invariant vectors,
and the space of $\fp^+_1$-invariant vectors is equal to
\begin{gather*} \cH_\tau(D,V)_{\tilde{K}_1}^{\fp^+_1}=\cH_\tau(D,V)\cap\cP(\fp^+_2,V). \end{gather*}
Thus if we write the decomposition of the above space under $\tilde{K}_1^\BC$ as
\begin{gather*} \cH_\tau(D,V)\cap\cP(\fp^+_2,V)\simeq \bigoplus_i m(\rho_i)(\rho_i,W_i), \end{gather*}
then $\cH_\tau(D,V)$ is decomposed under $\tilde{G}_1$ abstractly as
\begin{gather*}
\cH_\tau(D,V)_{\tilde{K}}|_{(\fg_1,\tilde{K}_1)} \simeq \bigoplus_i m(\rho_i)\cH_{\rho_i}(D_1,W_i)_{\tilde{K}_1},\\
\cH_\tau(D,V)|_{\tilde{G}_1} \simeq \hsum_i m(\rho_i)\cH_{\rho_i}(D_1,W_i)
\end{gather*}
(see \cite{JV}, \cite[Section 8]{Kmf1}, \cite{Ma}). Especially, if $\cH_\tau(D,V)_{\tilde{K}}=\cP(\fp^+, V)$, the decomposition of
$\cH_\tau(D,V)$ under $\tilde{G}_1$ corresponds to the decomposition of $\cP(\fp^+_2,V)\simeq\cP(\fp^+_2)\otimes V$ under $\tilde{K}_1^\BC$.
Now let $(\rho,W)$ be an irreducible representation of $\tilde{K}_1^\BC$ which appears in $\cP(\fp^+_2)\otimes V|_{\tilde{K}_1^\BC}$,
with $\tilde{K}_1$-invariant inner product $(u,v)_\rho$.
Our aim is to construct $\tilde{G}_1$-intertwining operators between $\cH_\tau(D,V)|_{\tilde{G}_1}$ and $\cH_\rho(D_1,W)$
explicitly. To do this, we assume $\cH_\tau(D,V)_{\tilde{K}}=\cP(\fp^+,V)$.

\subsection{Integral expression}

First we find the intertwining operators in integral expression.
Let $(\rho,W)$ be an irreducible representation of $\tilde{K}_1^\BC$ which appears in $\cP(\fp^+_2)\otimes V|_{\tilde{K}_1^\BC}$,
and let $\cF_{\tau\rho}^*\colon \cH_\tau(D,V)\to \cH_{\rho}(D_1,W)$ be a $\tilde{G}_1$-intertwining operator.
Then for any $y_1\in D_1$, the linear map $\cH_\tau(D,V)\to W$, $f\mapsto (\cF_{\tau\rho}^* f)(y_1)$ is continuous,
and by the Riesz representation theorem, there exists $\hat{\rK}_{y_1}\in\cH_\tau(D,V)\otimes \overline{W}$ such that
\begin{gather*} \langle f,\hat{\rK}_{y_1}\rangle_{\hat{\tau}}=(\cF_{\tau\rho}^* f)(y_1), \qquad f\in\cH_\tau(D,V),\quad y_1\in D_1. \end{gather*}
We write $\hat{\rK}(x;y_1):=\hat{\rK}_{y_1}(x)$ for $x\in D$, $y_1\in D_1$.
We identify $V\otimes \overline{W}$ and $\Hom(W,V)$ via the inner product of $W$.
Then $\cF_{\tau\rho}^*\colon \cH_\tau(D,V)\to \cH_{\rho}(D_1,W)$ and its adjoint operator $\cF_{\tau\rho}\colon \cH_\rho(D_1,W)\to \cH_\tau(D,V)$ are given by
\begin{alignat*}{3}
&(\cF_{\tau\rho}^* f)(y_1)=\langle f,\hat{\rK}(\cdot,y_1)\rangle_{\hat{\tau}},\qquad && f\in\cH_\tau(D,V),\quad y_1\in D_1,\\
& (\cF_{\tau\rho} f)(x)=\langle f,\hat{\rK}(x,\cdot)^*\rangle_{\hat{\rho}}, \qquad && f\in\cH_\rho(D,W),\quad x\in D.
\end{alignat*}
By the intertwining property, $\hat{\rK}(x;y_1)$ must satisfy
\begin{gather}\label{G1-invariance}
\hat{\rK}(gx;gy_1)=\tau(\kappa(g,x))\hat{\rK}(x;y_1)\rho(\kappa(g,y_1))^*
\end{gather}
for any $g\in\tilde{G}_1$, where $\kappa(g,x)$ is as (\ref{cocycle}). Thus we find a kernel function satisfying (\ref{G1-invariance}).

Let $\rK(x_2)\in\cP(\fp^+_2)\otimes V\otimes\overline{W}\simeq\cP(\fp^+_2,\Hom(W,V))$ be an operator-valued polynomial satisfying
\begin{gather}\label{K-invariance}
\rK(kx_2)=\tau(k)\rK(x_2)\rho(k)^{-1}, \qquad x_2\in \fp^+_2,\quad k\in \tilde{K}^\BC_1.
\end{gather}
Let $\Proj_2\colon \fp^+\to\fp^+_2$ be the orthogonal projection, and we define an operator-valued function $\hat{\rK}\in\cO(D\times \overline{D_1}, \Hom(W,V))$ by
\begin{gather}\label{Khat}
\hat{\rK}(x;y_1)=\hat{\rK}(x_1,x_2;y_1):=\tau(B(x,y_1))\rK(\Proj_2(x^{y_1})),
\end{gather}
where $x=(x_1,x_2)\in D$, $y_1\in D_1$. Then the following holds.
\begin{Proposition}\label{kernelfunc}
For any $x\in D$, $y_1\in D_1$, and $g\in\tilde{G}_1$, $\hat{\rK}(x;y_1)$ satisfies the identity \eqref{G1-invariance}.
\end{Proposition}
\begin{proof}
By (\ref{Bergman_property}), we have
\begin{gather*} \tau(B(gx,gy_1))=\tau(\kappa(g,x))\tau(B(x,y_1))\tau(\kappa(g,y_1))^*. \end{gather*}
Thus it suffices to show
\begin{gather*} \rK(\Proj_2((gx)^{gy_1}))=\tau(\kappa(g,y_1))^{*-1}\rK(\Proj_2(x^{y_1}))\rho(\kappa(g,y_1))^*. \end{gather*}
By $\tilde{K}^\BC_1$-invariance of $\rK(\cdot)$, this is equivalent to
\begin{gather*} \Proj_2((gx)^{gy_1})=\kappa(g,y_1)^{*-1}\Proj_2(x^{y_1}), \qquad x\in D,\quad y_1\in D_1,\quad g\in G_1. \end{gather*}
First we show
\begin{gather*} \Proj_2\big((gx)^{(\hat{\vartheta}g)y_1}\big)=\kappa(\hat{\vartheta}g,y_1)^{*-1}\Proj_2(x^{y_1}), \qquad x\in \fp^+,\quad y_1\in \fp^+_1 \end{gather*}
for $g=k\in K_1^\BC$ or $g=\exp(-z_1)$, $g=\exp(\vartheta w_1)\in G_1^\BC$ with $z_1,w_1\in\fp^+_1$,
when one side is well-defined, that is, we show
\begin{gather*}
\Proj_2\big((kx)^{k^{*-1}y_1}\big)=k\Proj_2(x^{y_1}), \\
\Proj_2\big((x-z_1)^{(y_1^{z_1})}\big)=B(z_1,y_1)\Proj_2(x^{y_1}), \\
\Proj_2\big((x^{w_1})^{y_1-w_1}\big)=\Proj_2(x^{y_1}).
\end{gather*}
In fact, the 1st and 3rd formulas follow from (\ref{quasiinv_k}), (\ref{quasiinv_add}),
and the 2nd formula follows from $(x-z_1)^{(y_1^{z_1})}=B(z_1,y_1)x^{y_1}-B(z_1,y_1)z_1^{y_1}$,
which is a consequence of~(\ref{quasiinv_twice}), and that $B(z_1,y_1)z_1^{y_1}\in\fp^+_1$ is annihilated by $\Proj_2$.
Since any $g\in G_1$ is written as the form $g=\exp(\vartheta w_1)k\exp(-z_1)$ with $z_1,w_1\in D_1$ and $k\in K^\BC_1$
(which is proved by using the $KAK$-decomposition and \cite[Part~III, Lemma~III.2.4]{FKKLR}),
the proposition follows from the cocycle condition of $\kappa$.
\end{proof}

We write $\cP(\fp^+_2,\Hom(W,V))^{\tilde{K}^\BC_1}:=\{\rK(x_2)\in \cP(\fp^+_2,\Hom(W,V))\colon \text{satisfying }(\ref{K-invariance})\}$,
$\cO(D\times \overline{D_1}$, $\Hom(W,V))^{\tilde{G}_1}
:=\{\hat{\rK}(x;y_1)\in\cO(D\times \overline{D_1}, \Hom(W,V))\colon \text{satisfying }(\ref{G1-invariance})\}$.
Then we have the following.
\begin{Lemma}
The linear map $\cO(D\times \overline{D_1}, \Hom(W,V))^{\tilde{G}_1}\to\cP(\fp^+_2,\Hom(W,V))^{\tilde{K}^\BC_1}$,
$\hat{\rK}(x;y_1)\mapsto \hat{\rK}(0,x_2;0)$ is bijective, and its inverse is given by $\rK(x_2)\mapsto \hat{\rK}(x;y_1)$ in \eqref{Khat}.
\end{Lemma}
\begin{proof}We write the restriction map by $\Rest$, that is, $(\Rest \hat{\rK})(x_2):=\hat{\rK}(0,x_2;0)$,
and write the linear map in~(\ref{Khat}) by $\Ext$.
For any $\hat{\rK}\in\cO(D\times \overline{D_1}, \Hom(W,V))^{\tilde{G}_1}$, clearly $(\Rest \hat{\rK})(x_2)$ satisfies
(\ref{K-invariance}), and therefore
\begin{gather*}
\big(\Rest \hat{\rK}\big)(x_2)\in \cO(D\cap \fp^+_2,\Hom(W,V))^{\tilde{K}^\BC_1}\\
\qquad {} \simeq \Hom_{\tilde{K}^\BC_1}(W,\cO(D\cap\fp^+_2,V)) \simeq\Hom_{\tilde{K}^\BC_1}(W,\cP(\fp^+_2,V))\simeq\cP(\fp^+_2,\Hom(W,V))^{\tilde{K}^\BC_1}
\end{gather*}
holds.
Also, for any $\rK(x_2)\in\cP(\fp^+_2,\allowbreak\Hom(W,V))^{\tilde{K}^\BC_1}$ we have
\begin{gather*} (\Rest\circ\Ext \rK)(x_2)=\tau(B(x_2,0))\rK\big(\Proj_2\big((x_2)^0\big)\big)=\rK(x_2). \end{gather*}
Thus $\Rest$ is surjective. Therefore it suffices to show that $\Rest$ is injective.
Suppose $\hat{\rK}_1$ and $\hat{\rK}_2\in\cO(D\times \overline{D_1}, \Hom(W,V))^{\tilde{G}_1}$ satisfy $\Rest\hat{\rK}_1=\Rest\hat{\rK}_2$.
Then by~(\ref{G1-invariance}),
\begin{gather*} \hat{\rK}_1(g.(0,x_2);g.0)=\hat{\rK}_2(g.(0,x_2);g.0) \end{gather*}
holds for any $g\in\tilde{G}_1$.
We fix $x_2\in D\cap \fp^+_2$, and set
\begin{gather*}
S_{x_2} :=\left\{ (g.(0,x_2);g.0)\in D\times\overline{D_1}\colon g\in G_1\right\}\subset D\times\overline{D_1}, \\
D_{x_2} :=\left\{ (g.(0,x_2);h.0)\in D\times\overline{D_1}\colon g,h\in G_1\right\}\subset D\times\overline{D_1}.
\end{gather*}
We show that $S_{x_2}$ contains a totally real submanifold of full dimension of $D_{x_2}$.
Let $\mathrm{pr}_1\colon D\times\overline{D_1}\to D\to D\cap\fp^+_1=D_1$, $\mathrm{pr}_2\colon D\times\overline{D_1}\to \overline{D_1}$
be the projections. Then since for every $x_2\in D\cap\fp^+_2$, $\{\exp(z).(0,x_2)\colon z\in\fp_1\}\subset D$ intersects transversally
with $\fp^+_2$ at $x_2$, the differential of $\mathrm{pr}_1|_{S_{x_2}}$ at $(0,x_2;0)$ is surjective.
Similarly, since $G_1$ acts transitively on $D_1$, the differential of $\mathrm{pr}_2|_{S_{x_2}}$ at $(0,x_2;0)$ is also surjective.
Therefore, $\mathrm{pr}_1|_{S_{x_2}}$ and $\mathrm{pr}_2|_{S_{x_2}}$ are both submersive near $(0,x_2;0)\in S_{x_2}$,
and $T_{(x;y_1)}S_{x_2}+JT_{(x;y_1)}S_{x_2}=T_{(x;y_1)}D_{x_2}$ holds on this neighborhood,
where $J$ is the complex structure of $D_{x_2}$. Hence $S_{x_2}$ contains a totally real submanifold of full dimension of $D_{x_2}$,
and $\hat{\rK}_1=\hat{\rK}_2$ holds on $D_{x_2}$ for each $x_2\in D\cap\fp^+_2$.
Finally, since $\bigcup_{x_2\in D\cap\fp^+_2} D_{x_2}$ contains an open subset of $D\times\overline{D_1}$,
$\hat{\rK}_1=\hat{\rK}_2$ holds on whole $D\times\overline{D_1}$. Therefore $\Rest$ is injective.
\end{proof}

Therefore if we assume $\cH_\tau(D,V)_{\tilde{K}}=\cP(\fp^+,V)$, then for any $(\rho,W)\subset \cP(\fp^+_2,V)$,
\begin{gather*}
 \overline{\Hom_{\tilde{G}_1}(\cH_\tau(D,V),\cH_\rho(D_1,W))}\simeq \Hom_{\tilde{G}_1}(\cH_\rho(D_1,W),\cH_\tau(D,V)) \\
\hphantom{\overline{\Hom_{\tilde{G}_1}(\cH_\tau(D,V),\cH_\rho(D_1,W))}}{} \simeq \cP(\fp^+_2,\Hom(W,V))^{\tilde{K}^\BC_1}\simeq \cO(D\times \overline{D_1}, \Hom(W,V))^{\tilde{G}_1}
\end{gather*}
holds, and by the above argument,
for any $\rK\in\cP(\fp^+_2,\Hom(W,V))^{\tilde{K}^\BC_1}$, $\hat{\rK}(x;y_1)$ in (\ref{Khat}) becomes
the kernel function of the intertwining operator.
Especially, $\hat{\rK}(\cdot;y_1)\in\cH_\tau(D,V)\otimes \overline{W}$ holds for any $y_1\in D_1$,
and $\hat{\rK}(x;\cdot)^*\in\cH_{\rho}(D_1,W)\otimes \overline{V}$ holds for any $x\in D$.
Also, even if $\cH_\tau(D,V)_{\tilde{K}}\ne\cP(\fp^+)\otimes V$, since $\hat{\rK}(x;\cdot)$ is a bounded function on $D_1$ for any $x\in D$,
$\hat{\rK}(x;\cdot)^*\in\cH_{\rho}(D_1,W)\otimes \overline{V}$ holds
if we assume $\cH_\rho(D_1,W)$ is a holomorphic discrete series representation.
\begin{Corollary}\label{integral_expression}
Let $(\tau,V)$ be a $\tilde{K}^\BC$-module, and $(\rho,W)$ be a $\tilde{K}_1^\BC$-module
which appears in $\cP(\fp^+_2)\otimes V|_{\tilde{K}_1^\BC}$.
Let $\rK(x_2)\in \cP(\fp^+_2)\otimes V\otimes\overline{W}\simeq\cP(\fp^+_2,\Hom(W,V))$ be an operator-valued polynomial
satisfying \eqref{K-invariance}, and define $\hat{\rK}(x;y_1)\in\cO(D\times\overline{D_1},\Hom(W,V))$ by \eqref{Khat}.
\begin{enumerate}\itemsep=0pt
\item[$(1)$] Assume $\cH_\tau(D,V)_{\tilde{K}}=\cP(\fp^+,V)$. Then the linear maps
\begin{gather*}
\cF_{\tau\rho}^*\colon \ \cH_\tau(D,V)\to\cH_\rho(D_1,W), \\
(\cF_{\tau\rho}^*f)(y_1):=\big\langle f,\hat{\rK}(\cdot;y_1)\big\rangle_{\hat{\tau}}
=C_\tau\int_D\hat{\rK}(x;y_1)^*\tau\big(B(x)^{-1}\big)f(x)h(x)^{-p}{\rm d}x, \\
\cF_{\tau\rho}\colon \ \cH_\rho(D_1,W)\to \cH_\tau(D,V), \\
(\cF_{\tau\rho}f)(x):=\big\langle f,\hat{\rK}(x;\cdot)^*\big\rangle_{\hat{\rho}}
=C_\rho\int_{D_1}\hat{\rK}(x;y_1)\rho\big(B(y_1)^{-1}\big)f(y_1)h_1(y_1)^{-p_1}{\rm d}y_1
\end{gather*}
intertwine the $\tilde{G}_1$-action. Here the second equalities hold only if $\cH_\tau(D{,}V)$ resp.~$\cH_\rho(D_1{,}W)$ are holomorphic discrete series representations.
\item[$(2)$] Assume $\cH_\rho(D_1,W)$ is a holomorphic discrete series representation.
Then the linear map $\cF_{\tau\rho}\colon \cH_\rho(D_1,W)\to \cO_\tau(D,V)$,
\begin{gather*} (\cF_{\tau\rho}f)(x):=\big\langle f,\hat{\rK}(x;\cdot)^*\big\rangle_{\hat{\rho}}
=C_\rho\int_{D_1}\hat{\rK}(x;y_1)\rho\big(B(y_1)^{-1}\big)f(y_1)h_1(y_1)^{-p_1}{\rm d}y_1 \end{gather*}
intertwines the $\tilde{G}_1$-action.
\end{enumerate}
\end{Corollary}
These operators are defined as maps from $\cH_\tau(D,V)$ resp.~$\cH_\rho(D_1,W)$, but in fact these extend continuously to maps between
$\cO_\tau(D,V)$ and $\cO_\rho(D_1,W)$ if $\cH_\tau(D,V)$ or $\cH_\rho(D_1,W)$ is a~holomorphic discrete series representation.
\begin{Theorem}\label{conti_ext}\quad
\begin{enumerate}\itemsep=0pt
\item[$(1)$] Assume $\cH_\tau(D,V)$ is a holomorphic discrete series representation. Then
the linear map $\cF_{\tau\rho}^*\colon\! \cH_\tau(D,V)\!\to\!\cH_\rho(D_1{,}W)$ extends continuously to the map $\cF_{\tau\rho}^*\colon\! \cO_\tau(D,V)\!\to\!\cO_\rho(D_1{,}W)$.
\item[$(2)$] Assume $\cH_\rho(D_1,W)$ is a holomorphic discrete series representation. Then
the linear map $\cF_{\tau\rho}\colon\! \cH_\rho(D_1,W)\!\to\! \cO_\tau(D{,}V)$ extends continuously to the map $\cF_{\tau\rho}\colon\! \cO_\rho(D_1,W)\!\to \! \cO_\tau(D{,}V)$.
\end{enumerate}
\end{Theorem}
\begin{proof}
We only prove (2). By the $\tilde{G}_1$-invariance of $\hat{\rK}$, for $k\in \tilde{K}_1$ we have
\begin{gather*} (\cF_{\tau\rho}f)(x)=\big\langle f(y_1),\hat{\rK}(x;y_1)^*\big\rangle_{\hat{\rho},y_1}
=\big\langle \rho(k)^{-1}f(ky_1),\hat{\rK}\big(k^{-1}x;y_1\big)^*\tau(k)^*\big\rangle_{\hat{\rho},y_1}. \end{gather*}
Let $\hbar=-\sqrt{-1}z\in\sqrt{-1}\fz(\fk)$ be the element such that $\operatorname{ad}(\hbar)|_{\fp^+}=I_{\fp^+}$.
Then for $t\in\sqrt{-1}\BR$ and $v\in V$, by setting $k={\rm e}^{-t\hbar}$ we have
\begin{gather*} ((\cF_{\tau\rho}f)(x),v)_\tau=\chi_\rho\big({\rm e}^{t\hbar}\big)\chi_\tau\big({\rm e}^{-t\hbar}\big)\big\langle f\big({\rm e}^{-t}y_1\big),
\hat{\rK}\big({\rm e}^tx;y_1\big)^*v\big\rangle_{\hat{\rho},y_1}, \end{gather*}
where we write $\rho({\rm e}^{t\hbar})=\chi_\rho({\rm e}^{t\hbar})I_W$, $\tau({\rm e}^{t\hbar})=\chi_\tau({\rm e}^{t\hbar})I_V$.
Here, though $\chi_\rho$ and $\chi_\tau$ are defined as a~function on $Z(\tilde{K})\subset \tilde{G}$,
their ratio $\chi_\rho\chi_\tau^{-1}$ is well-defined as a function on $Z(K)\subset G$.
By analytic continuation this holds for $t\in\BC$ such that $|{\rm e}^t|\ge 1$ and $|{\rm e}^t x|_\infty <1$.
Then for $|x|_\infty<{\rm e}^{-t}\le 1$ we have
\begin{align*}
|((\cF_{\tau\rho}f)(x),v)_\tau|&=\big|\chi_\rho({\rm e}^{t\hbar})\chi_\tau({\rm e}^{-t\hbar})\big\langle f\big({\rm e}^{-t}y_1\big),
\hat{\rK}\big({\rm e}^tx;y_1\big)^*v\big\rangle_{\hat{\rho},y_1}\big| \\
&\le \chi_\rho({\rm e}^{t\hbar})\chi_\tau\big({\rm e}^{-t\hbar}\big)\big\Vert f\big({\rm e}^{-t}y_1\big)\big\Vert_{\hat{\rho},y_1}
\big\Vert\hat{\rK}\big({\rm e}^tx;y_1\big)^*v\big\Vert_{\hat{\rho},y_1}.
\end{align*}
Now since $\cH_\rho(D_1,W)$ is a holomorphic discrete series representation, we have
\begin{align*}
\big\Vert f\big({\rm e}^{-t}y_1\big)\big\Vert_{\hat{\rho},y_1}^2
&=C_\rho\int_{D_1}\big(\rho(B(y_1)^{-1})f\big({\rm e}^{-t}y_1\big),f\big({\rm e}^{-t}y_1\big)\big)_\rho h_1(y_1)^{-p_1}{\rm d}y_1 \\
&\le C_\rho\int_{D_1}\big|\rho\big(B(y_1)^{-1}\big)\big|_{\rho,\mathrm{op}}h_1(y_1)^{-p_1}{\rm d}y_1 \sup_{y_1\in D_1}\big|f\big({\rm e}^{-t}y_1\big)\big|_\rho^2 \\
&=C_\rho^{\prime 2}\sup_{|y_1|_\infty\le {\rm e}^{-t}} |f(y_1)|_\rho^2.
\end{align*}
Therefore for $|x|_\infty<{\rm e}^{-t}$ we have
\begin{gather*}
 |((\cF_{\tau\rho}f)(x),v)_\tau|
\le C_\rho'\chi_\rho\big({\rm e}^{t\hbar}\big)\chi_\tau\big({\rm e}^{-t\hbar}\big)\sup_{|y_1|_\infty\le {\rm e}^{-t}}|f(y_1)|_\rho
\big(\big\langle\hat{\rK}\big({\rm e}^tx;\cdot\big)^*,\hat{\rK}\big({\rm e}^tx;\cdot\big)^*\big\rangle_{\hat{\rho}}v,v\big)_\tau^{1/2},
\end{gather*}
and since $\big\langle\hat{\rK}(x;\cdot)^*,\hat{\rK}(x;\cdot)^*\big\rangle_{\hat{\rho}}\in C^\omega(D,\End(V))$,
$\cF_{\tau\rho}$ extends continuously to the map from $\cO_\rho(D_1,W)$ to $\cO_\tau(D,V)$.
\end{proof}

\subsection{Differential expression}
Next we will establish differential expressions for the intertwining operators.
To do this for $\cF_{\tau\rho}\colon \cH_\rho(D_1,W)\to \cH_\tau(D,V)$, we assume that $(G,G_1)$ is a symmetric pair, that is,
there exists an involution $\sigma$ of $G$ (without loss of generality we assume $\sigma\vartheta=\vartheta\sigma$) such that $G_1=(G^\sigma)_0$.
A~symmetric pair which satisfies the assumption (\ref{assumption}) is called a symmetric pair of \textit{holomorphic type}.
We extend $\sigma$ on $\fg^\BC$ holomorphically. Then this defines the involution on $\fp^+$ and also on~$D$.

Since the reproducing kernel of $\cP(\fp^+,V)$ with respect to the Fischer norm
is given by ${\rm e}^{(x|z)_{\fp^+}}I_V$, that is,
\begin{gather*} f(x)=\frac{1}{\pi^n}\int_{\fp^+}{\rm e}^{(x|z)_{\fp^+}}f(z){\rm e}^{-|z|_{\fp^+}^2}{\rm d}z \end{gather*}
holds for $f\in \cH_\tau(D,V)_{\tilde{K}}=\cP(\fp^+,V)$, we have
\begin{align*}
(\cF_{\tau\rho}^*f)(y_1)&=\left\langle \frac{1}{\pi^n}\int_{\fp^+}{\rm e}^{(\cdot|z)_{\fp^+}}f(z){\rm e}^{-|z|_{\fp^+}^2}{\rm d}z,
\hat{\rK}(\cdot;y_1)\right\rangle_{\hat{\tau}} \\
&=\frac{1}{\pi^n}\int_{\fp^+}\big\langle {\rm e}^{(\cdot|z)_{\fp^+}}I_V,\hat{\rK}(\cdot;y_1)\big\rangle_{\hat{\tau}}f(z){\rm e}^{-|z|_{\fp^+}^2}{\rm d}z \\
&=\frac{1}{\pi^n}\int_{\fp^+}\cF_{\tau\rho}^*\big({\rm e}^{(\cdot|z)_{\fp^+}}I_V\big)(y_1)f(z){\rm e}^{-|z|_{\fp^+}^2}{\rm d}z,
\end{align*}
and similarly, for $f\in \cH_\rho(D_1,W)_{\tilde{K}_1}=\cP(\fp^+_1,W)$ we have
\begin{align*}
(\cF_{\tau\rho}f)(x)&=\left\langle \frac{1}{\pi^{n_1}}\int_{\fp^+_1}{\rm e}^{(\cdot|w_1)_{\fp^+_1}}f(w_1){\rm e}^{-|w_1|_{\fp^+_1}^2}{\rm d}w_1,
\hat{\rK}(x;\cdot)^*\right\rangle_{\hat{\rho}} \\
&=\frac{1}{\pi^{n_1}}\int_{\fp^+_1}\big\langle {\rm e}^{(\cdot|w_1)_{\fp^+_1}}I_W,\hat{\rK}(x;\cdot)^*\big\rangle_{\hat{\rho}}
f(w_1){\rm e}^{-|w_1|_{\fp^+_1}^2}{\rm d}w_1 \\
&=\frac{1}{\pi^{n_1}}\int_{\fp^+_1}\cF_{\tau\rho}\big({\rm e}^{(\cdot|w_1)_{\fp^+_1}}I_W\big)(x)f(w_1){\rm e}^{-|w_1|_{\fp^+_1}^2}{\rm d}w_1.
\end{align*}
Now we have
\begin{Lemma}\quad
\begin{enumerate}\itemsep=0pt
\item[$(1)$] $\ds \cF_{\tau\rho}^*\big({\rm e}^{(\cdot|z)_{\fp^+}}I_V\big)(y_1)=\cF_{\tau\rho}^*\big({\rm e}^{(\cdot|z)_{\fp^+}}I_V\big)(0){\rm e}^{(y_1|z)_{\fp^+}}$.
\item[$(2)$] If $(G,G_1)$ is symmetric, then
\begin{gather*} \cF_{\tau\rho}\big({\rm e}^{(\cdot|w_1)_{\fp^+_1}}I_W\big)(x)=\cF_{\tau\rho}\big({\rm e}^{(\cdot|w_1)_{\fp^+_1}}I_W\big)(0,x_2){\rm e}^{(x_1|w_1)_{\fp^+_1}}. \end{gather*}
\end{enumerate}
\end{Lemma}
\begin{proof}
(1) Since $\cF_{\tau\rho}^*$ intertwines the $\tilde{G}_1$-action, it also intertwines the $\fg^\BC_1$-action.
Especially, since $\fp^+_1\subset \fg^\BC_1$ acts as a 1st-order differential operator with constant coefficients, we have
\begin{align*}
\frac{{\rm d}}{{\rm d}t}\cF_{\tau\rho}^*\big({\rm e}^{(\cdot|z)_{\fp^+}}I_V\big)(ty_1)
&=\left.\frac{{\rm d}}{{\rm d}s}\right|_{s=0}\cF_{\tau\rho}^*\big({\rm e}^{(\cdot|z)_{\fp^+}}I_V\big)(ty_1+sy_1) \\
&=\left.\frac{{\rm d}}{{\rm d}s}\right|_{s=0}\cF_{\tau\rho}^*\big({\rm e}^{(\cdot+sy_1|z)_{\fp^+}}I_V\big)(ty_1)\\
& =\left.\frac{{\rm d}}{{\rm d}s}\right|_{s=0}\cF_{\tau\rho}^*\big({\rm e}^{(\cdot|z)_{\fp^+}}I_V\big)(ty_1){\rm e}^{s(y_1|z)_{\fp^+}} \\
&=\cF_{\tau\rho}^*\big({\rm e}^{(\cdot|z)_{\fp^+}}I_V\big)(ty_1)(y_1|z)_{\fp^+}.
\end{align*}
Therefore, as functions of $t$, both
\begin{gather*} \cF_{\tau\rho}^*\big({\rm e}^{(\cdot|z)_{\fp^+}}I_V\big)(ty_1)\qquad \text{and} \qquad
\cF_{\tau\rho}^*\big({\rm e}^{(\cdot|z)_{\fp^+}}I_V\big)(0){\rm e}^{t(y_1|z)_{\fp^+}} \end{gather*}
satisfy the same differential equation with the same initial condition, and thus they coincide.
By substituting $t=1$, we get the desired formula.

(2) First we note that under the assumption that $(G,G_1)$ is symmetric,
if $x=(x_1,x_2)\in D$, where $x_1\in\fp^+_1$, $x_2\in\fp^+_2$, then $(tx_1,x_2)\in D$ holds for any $t\in[-1,1]$,
because $-\sigma(x_1,x_2)=(-x_1,x_2)\in D$ and $D$ is convex.
Therefore $\cF_{\tau\rho}\big({\rm e}^{(\cdot|w_1)_{\fp^+_1}}I_W\big)(tx_1,x_2)$ is well-defined for any $t\in[-1,1]$ if $(G,G_1)$ is symmetric.
Then as for (1), we can show that $\cF_{\tau\rho}\big({\rm e}^{(\cdot|w_1)_{\fp^+_1}}I_W\big)(tx_1,x_2)$ and
$\cF_{\tau\rho}\big({\rm e}^{(\cdot|w_1)_{\fp^+_1}}I_W\big)(0,x_2){\rm e}^{t(x_1|w_1)_{\fp^+_1}}$ satisfy the same differential equation
with the same initial condition, and thus they coincide.
\end{proof}

Thus we define $F_{\tau\rho}^*(z)\in\cP(\overline{\fp^+},\Hom(V,W))$ and
$F_{\tau\rho}(x_2;w_1)\in\cO((D\cap\fp^+_2)\times\overline{\fp^+_1},\allowbreak\Hom(W,V))$ by
\begin{gather*}
F_{\tau\rho}^*(z)=F_{\tau\rho}^*(z_1,z_2):=\cF_{\tau\rho}^*\big({\rm e}^{(\cdot|z)_{\fp^+}}I_V\big)(0)
=\big\langle {\rm e}^{(\cdot|z)_{\fp^+}}I_V,\hat{\rK}(\cdot;0)\big\rangle_{\hat{\tau}} \\
\hphantom{F_{\tau\rho}^*(z)}{} =\big\langle {\rm e}^{(\cdot|z)_{\fp^+}}I_V,\rK(\Proj_2(\cdot))\big\rangle_{\hat{\tau}}
=C_\tau\int_D \rK(x_2)^*\tau\big(B(x)^{-1}\big){\rm e}^{(x|z)_{\fp^+}}h(x)^{-p}{\rm d}x,\\
F_{\tau\rho}(x_2;w_1)=\cF_{\tau\rho}\big({\rm e}^{(\cdot|w_1)_{\fp^+_1}}I_W\big)(0,x_2)
=\big\langle {\rm e}^{(y_1|w_1)_{\fp^+_1}}I_W,\hat{\rK}(0,x_2;y_1)^*\big\rangle_{\hat{\rho},y_1}\\
\hphantom{F_{\tau\rho}(x_2;w_1)}{}=\big\langle {\rm e}^{(y_1|w_1)_{\fp^+_1}}I_W,\bigl(\tau(B(x_2,y_1))\rK(\Proj_2((x_2)^{y_1}))\bigr)^*\big\rangle_{\hat{\rho},y_1}\\
\hphantom{F_{\tau\rho}(x_2;w_1)}{}=C_\rho\int_{D_1} \tau(B(x_2,y_1))\rK(\Proj_2((x_2)^{y_1}))\rho\big(B(y_1)^{-1}\big){\rm e}^{(y_1|w_1)_{\fp^+_1}}h_1(y_1)^{-p}{\rm d}y_1
\end{gather*}
(where the last equalities hold only if $\cH_\tau(D,V)$ resp.~$\cH_\rho(D_1,W)$ is a holomorphic discrete series representation).
Then we have
\begin{align*}
(\cF_{\tau\rho}^*f)(y_1)&=\frac{1}{\pi^n}\int_{\fp^+}F_{\tau\rho}^*(z_1,z_2){\rm e}^{(y_1|z)_{\fp^+}}f(z){\rm e}^{-|z|_{\fp^+}^2}{\rm d}z\\
&=\frac{1}{\pi^n}\left.\int_{\fp^+}F_{\tau\rho}^*(z_1,z_2){\rm e}^{(x|z)_{\fp^+}}f(z){\rm e}^{-|z|_{\fp^+}^2}{\rm d}z\right|_{x_1=y_1,x_2=0}\\
&=F_{\tau\rho}^*\left(\left.\overline{\frac{\partial}{\partial x_1}}\right|_{x_1=y_1},
\left.\overline{\frac{\partial}{\partial x_2}}\right|_{x_2=0}\right)
\frac{1}{\pi^n}\int_{\fp^+}{\rm e}^{(x|z)_{\fp^+}}f(z){\rm e}^{-|z|_{\fp^+}^2}{\rm d}z\\
&=F_{\tau\rho}^*\left.\left(\overline{\frac{\partial}{\partial x_1}},\overline{\frac{\partial}{\partial x_2}}\right)\right|_{x_1=y_1,x_2=0}f(x),
\end{align*}
and similarly we have
\begin{gather*}
(\cF_{\tau\rho}f)(x) =\frac{1}{\pi^{n_1}}\int_{\fp^+_1}F_{\tau\rho}(x_2;w_1){\rm e}^{(x_1|w_1)_{\fp^+_1}}f(w_1){\rm e}^{-|w_1|_{\fp^+_1}^2}{\rm d}w_1
=F_{\tau\rho}\left.\left(x_2;\overline{\frac{\partial}{\partial y_1}}\right)\right|_{y_1=x_1}f(y_1).
\end{gather*}
Here, for anti-holomorphic polynomial $f\in\cP(\overline{\fp^+})$, let $f\big(\overline{\frac{\partial}{\partial x}}\big)$ be the
holomorphic differential operator characterized by
\begin{gather*} f\left(\overline{\frac{\partial}{\partial x}}\right){\rm e}^{(x|y)_{\fp^+}}=f(y), \end{gather*}
and similarly for anti-holomorphic function $f\in\cO(\overline{\fp^+_1})$, let $f\big(\overline{\frac{\partial}{\partial x_1}}\big)$ be the
operator characterized by
\begin{gather*} f\left(\overline{\frac{\partial}{\partial x_1}}\right){\rm e}^{(x_1|y_1)_{\fp^+_1}}=f(y_1). \end{gather*}

Next we describe $\Proj_2((x_2)^{y_1})$, and $\tau(B(x_2,y_1))$ when $\tau=\chi^{-\lambda}$ is one-dimensional, where $\chi$ is as (\ref{char}),
namely $\tau(B(x_2,y_1))=\chi^{-\lambda}(B(x_2,y_1))=h(x_2,y_1)^{-\lambda}$.
\begin{Proposition}\label{projprop}
Assume $(G,G_1)$ is a symmetric pair.
\begin{enumerate}\itemsep=0pt
\item[$(1)$] $\Proj_2((x_2)^{y_1})=(x_2)^{Q(y_1)x_2}$.
\item[$(2)$] Assume $G$ is simple. Then $h(x_2,y_1)^2=h(Q(x_2)y_1,y_1)=h(x_2,Q(y_1)x_2)$.
\end{enumerate}
\end{Proposition}
\begin{Lemma}\label{projlemma}
Let $x,y\in\fp^+$.
\begin{enumerate}\itemsep=0pt
\item[$(1)$] $B(-x,y)B(x,y)=B(Q(x)y,y)=B(x,Q(y)x)$.
\item[$(2)$] $x^y=(Q(x)y)^y+x^{Q(y)x}$.
\end{enumerate}
\end{Lemma}
\begin{proof}
(1) Use \cite[Part V, Proposition I.5.1, (J4.2), (J4.2$'$)]{FKKLR} and $B(x,-y)=B(-x,y)$, $Q(-x)=Q(x)$.

(2) Both sides are computed as
\begin{align*}
({\rm l.h.s.})&=B(x,y)^{-1}(x-Q(x)y),\\
({\rm r.h.s.})&=B(Q(x)y,y)^{-1}(Q(x)y-Q(Q(x)y)y)+B(x,Q(y)x)^{-1}(x-Q(x)Q(y)x) \\
&=B(x,y)^{-1}B(-x,y)^{-1}(x+Q(x)y-Q(x)Q(y)x-Q(x)Q(y)Q(x)y),
\end{align*}
where we used (1) and \cite[Part V, Proposition I.4.1]{FKKLR}. Thus it suffices to show
\begin{gather*} B(-x,y)(x-Q(x)y)=x+Q(x)y-Q(x)Q(y)x-Q(x)Q(y)Q(x)y. \end{gather*}
In fact, we have
\begin{gather*}
B(-x,y)(x-Q(x)y) =(I+D(x,y)+Q(x)Q(y))(x-Q(x)y) \\
\qquad{}=x+D(x,y)x+Q(x)Q(y)x-Q(x)y-D(x,y)Q(x)y-Q(x)Q(y)Q(x)y \\
\qquad{}=x+D(x,y)x+Q(x)Q(y)x-Q(x)y-Q(x)D(y,x)y-Q(x)Q(y)Q(x)y \\
\qquad{}=x+2Q(x)y+Q(x)Q(y)x-Q(x)y-2Q(x)Q(y)x-Q(x)Q(y)Q(x)y \\
\qquad{}=x+Q(x)y-Q(x)Q(y)x-Q(x)Q(y)Q(x)y,
\end{gather*}
where we used \cite[Part V, Proposition I.2.1(J1)]{FKKLR} at the 3rd equality, and $D(x,y)x=2Q(x)y$ at the 4th equality.
Thus the lemma follows.
\end{proof}

\begin{proof}[Proof of Proposition \ref{projprop}]
(1) When $x_2\in\fp^+_2$, $y_1\in\fp^+_1$, we have $(Q(x_2)y_1)^{y_1}\in\fp^+_1$, $(x_2)^{Q(y_1)x_2}\allowbreak\in\fp^+_2$. Therefore
\begin{gather*} \Proj_2((x_2)^{y_1})=\Proj_2\big((Q(x_2)y_1)^{y_1}+(x_2)^{Q(y_1)x_2}\big)=(x_2)^{Q(y_1)x_2}. \end{gather*}

(2) We extend $\sigma$ on $\fg^\BC$ holomorphically. Then since $\sigma$ acts by $+1$ on $\fp^+_1$ and $-1$ on $\fp^+_2$,
$B(-x_2,y_1)=\sigma B(x_2,y_1)\sigma$ holds. Therefore by Lemma \ref{projlemma},
\begin{gather*}
\sigma B(x_2,y_1)\sigma B(x_2,y_1)=B(Q(x_2)y_1,y_1)=B(x_2,Q(y_1)x_2), \\
\therefore \Det(B(x_2,y_1))^2=\Det(B(Q(x_2)y_1,y_1))=\Det(B(x_2,Q(y_1)x_2)).
\end{gather*}
Since $\Det(B(x_2,y_1))=h(x_2,y_1)^p$, the proposition follows.
\end{proof}

Therefore when $(G,G_1)$ is symmetric, we have
\begin{gather*}
 F_{\tau\rho}(x_2;w_1)
:=\big\langle {\rm e}^{(y_1|w_1)_{\fp^+_1}}I_W,\bigl(\tau(B(x_2,y_1))\rK\big((x_2)^{Q(y_1)x_2}\big) \bigr)^*\big\rangle_{\hat{\rho},y_1},
\end{gather*}
and moreover if $G$ is simple and $\tau=\chi^{-\lambda}$ is one-dimensional, we have
\begin{align*}
F_{\tau\rho}(x_2;w_1)&=\big\langle {\rm e}^{(y_1|w_1)_{\fp^+_1}}I_W,
\big(h(x_2,Q(y_1)x_2)^{-\lambda/2}\rK\big((x_2)^{Q(y_1)x_2}\big)\big)^*\big\rangle_{\hat{\rho},y_1}\\
&=\big\langle {\rm e}^{(y_1|w_1)_{\fp^+_1}}I_W,
\big(h(Q(x_2)y_1,y_1)^{-\lambda/2}\rK\big((x_2)^{Q(y_1)x_2}\big)\big)^*\big\rangle_{\hat{\rho},y_1}.
\end{align*}
We summarize the above results.
\begin{Theorem}\label{main}
Let $(\tau,V)$ be a $\tilde{K}^\BC$-module, and $(\rho,W)$ be a $\tilde{K}_1^\BC$-module
which appears in $\cP(\fp^+_2)\otimes V|_{\tilde{K}_1^\BC}$.
Let $\rK(x_2)\in \cP(\fp^+_2)\otimes V\otimes\overline{W}\simeq\cP(\fp^+_2,\Hom(W,V))$ be an operator-valued polynomial
satisfying \eqref{K-invariance}.
\begin{enumerate}\itemsep=0pt
\item[$(1)$] Assume $\cH_\tau(D,V)_{\tilde{K}}=\cP(\fp^+,V)$.
We define $F_{\tau\rho}^*(z)\in\cP(\overline{\fp^+},\Hom(V,W))$ by
\begin{gather*} F_{\tau\rho}^*(z)=F_{\tau\rho}^*(z_1,z_2):=\big\langle {\rm e}^{(\cdot|z)_{\fp^+}}I_V,\rK(\Proj_2(\cdot))\big\rangle_{\hat{\tau}}. \end{gather*}
Then the linear map
\begin{gather*}
\cF_{\tau\rho}^*\colon \cH_\tau(D,V)_{\tilde{K}}\to\cH_\rho(D_1,W)_{\tilde{K}_1}, \qquad
(\cF_{\tau\rho}^*f)(x_1)
=F_{\tau\rho}^*\left.\left(\overline{\frac{\partial}{\partial x_1}},\overline{\frac{\partial}{\partial x_2}}\right)\right|_{x_2=0}f(x)
\end{gather*}
intertwines the $(\fg_1,\tilde{K}_1)$-action.
\item[$(2)$] Assume $(G,G_1)$ is symmetric, and also assume ``$\cH_\tau(D,V)_{\tilde{K}}=\cP(\fp^+,V)$'' or ``$\cH_\rho(D_1,W)$ is
a holomorphic discrete series representation''.
We define $F_{\tau\rho}(x_2;w_1)\in\cO((D\cap\fp^+_2)\times\overline{\fp^+_1},$ $\Hom(W,V))$ by
\begin{gather*}
 F_{\tau\rho}(x_2;w_1)
:=\big\langle {\rm e}^{(y_1|w_1)_{\fp^+_1}}I_W,\bigl(\tau(B(x_2,y_1))\rK\big((x_2)^{Q(y_1)x_2}\big)\bigr)^*\big\rangle_{\hat{\rho},y_1}.
\end{gather*}
Then the linear map
\begin{gather*}
\cF_{\tau\rho}\colon \cH_\rho(D_1,W)_{\tilde{K}_1}\to \cO_\tau(D,V)_{\tilde{K}}, \qquad
(\cF_{\tau\rho}f)(x) =F_{\tau\rho}\left(x_2;\overline{\frac{\partial}{\partial x_1}}\right)f(x_1)
\end{gather*}
intertwines the $(\fg_1,\tilde{K}_1)$-action.
\end{enumerate}
\end{Theorem}

\begin{Remark}
For $w\in\fp^+$, we define the differential operator $\cB_\tau(w)$ on $\cP(\overline{\fp^+},V)$ by
\begin{gather*} \cB_\tau(w)f(z):=\sum_{\alpha\beta}\frac{1}{2}(Q(e_\alpha,e_\beta)w|z)_{\fp^+}
\frac{\partial^2f}{\partial\bar{z}_\alpha\partial\bar{z}_\beta}(z)
+\sum_{\alpha}{\rm d}\tau([e_\alpha,\vartheta w])\frac{\partial f}{\partial \bar{z}_\alpha}(z), \end{gather*}
where $\{e_\alpha\}$ is a basis of $\fp^+$, with the dual basis $\{e_\alpha^\vee\}$,
and $\frac{\partial}{\partial\bar{z}_\alpha}$ is the anti-holomorphic directional derivative along $e_\alpha^\vee$.
Then this is a generalization of the Bessel operator $\cB_\nu$ in~\cite{D} or \cite[Section~XV.2]{FK}.
Then for $w_1\in\fp^+_1$, $\cB_\tau(w_1)$ annihilates $F_{\tau\rho}^*(z)$, because
\begin{gather*}
 (\cB_\tau(w_1))_zF_{\tau\rho}^*(z)
=(\cB_\tau(w_1))_z\big\langle {\rm e}^{(x|z)_{\fp^+}}I_V,\rK(\Proj_2(x))\big\rangle_{\hat{\tau},x} \\
 =\big\langle ((Q(x)w_1|z)_{\fp^+}+{\rm d}\tau([x,\vartheta w_1])){\rm e}^{(x|z)_{\fp^+}}I_V,\rK(\Proj_2(x))\big\rangle_{\hat{\tau},x} \\
 =\big\langle {\rm d}\hat{\tau}(-\vartheta w_1)_x{\rm e}^{(x|z)_{\fp^+}}I_V,\rK(\Proj_2(x))\big\rangle_{\hat{\tau},x}
=\big\langle {\rm e}^{(x|z)_{\fp^+}}I_V,{\rm d}\hat{\tau}(w_1)_x\rK(\Proj_2(x))\big\rangle_{\hat{\tau},x} \\
 =\left\langle {\rm e}^{(x|z)_{\fp^+}}I_V,\left.\frac{{\rm d}}{{\rm d}t}\right|_{t=0}\rK(\Proj_2(x-tw_1))\right\rangle_{\hat{\tau},x}
=\left\langle {\rm e}^{(x|z)_{\fp^+}}I_V,\left.\frac{{\rm d}}{{\rm d}t}\right|_{t=0}\rK(\Proj_2(x))\right\rangle_{\hat{\tau},x}=0.
\end{gather*}
This differential equation coincides with $\widehat{{\rm d}\pi_\mu}$ on $\fn_+$ appeared in Proposition 3.10 or Section 4.4, Step 1
of \cite{KP1}, and thus the operator $\cF_{\tau\rho}^*$ coincides with the one given by the F-method.
\end{Remark}
These operators are defined as maps from the space of polynomials $\cH_\tau(D,V)_{\tilde{K}}$ resp.\ \linebreak $\cH_\rho(D_1,W)_{\tilde{K}_1}$,
but in fact these are well-defined as maps between
$\cO_\tau(D,V)$ and $\cO_\rho(D_1,W)$ in the following sense.
\begin{Theorem}\label{extend}\quad
\begin{enumerate}\itemsep=0pt
\item[$(1)$] $\cF_{\tau\rho}^*$ is well-defined as the map $\cF_{\tau\rho}^*\colon \cO_\tau(D,V)\to\cO_\rho(D_1,W)$.
\item[$(2)$] Assume $\cH_\rho(D_1\!{,}W)$ is a holomorphic discrete series representation. Then for  $f{\in} \cO_\rho(D_1\!{,}W)$,
$\cF_{\tau\rho}f(x)$ converges uniformly on every compact subset in $\{x=x_1+x_2\in D\colon |x_1|_\infty+|x_2|_\infty\allowbreak <1\}$,
and it continues holomorphically on whole $D$. Especially $\cF_{\tau\rho}$ is well-defined as the map
$\cF_{\tau\rho}\colon \cO_\rho(D_1,W)\to \cO_\tau(D,V)$.
\end{enumerate}
\end{Theorem}
\begin{proof}
(1) Clear since $\cF_{\tau\rho}^*$ is a finite-order differential operator.

(2) First we decompose $F_{\tau\rho}(x_2;w_1)$ as the sum of homogeneous polynomials in $w_1$.
\begin{align*}
F_{\tau\rho}(x_2;w_1) &=\big\langle {\rm e}^{(y_1|w_1)_{\fp^+_1}}I_W,\bigl(\tau(B(x_2,y_1))\rK(\Proj_2((x_2)^{y_1}))\bigr)^*\big\rangle_{\hat{\rho},y_1} \\
&=\sum_{n=0}^\infty \frac{1}{n!}\big\langle (y_1|w_1)_{\fp^+_1}^nI_W, \hat{\rK}(x_2;y_1)^*\big\rangle_{\hat{\rho},y_1} =:\sum_{n=0}^\infty F_n(x_2;w_1).
\end{align*}
Then $\cF_{\tau\rho}$ is written as
\begin{gather*} (\cF_{\tau\rho}f)(x)=\sum_{n=0}^\infty F_n\left(x_2;\overline{\frac{\partial}{\partial x_1}}\right)f(x_1). \end{gather*}
Now by the $\tilde{G}_1$-invariance of $\hat{\rK}$, for $k\in \tilde{K}_1$ we have
\begin{gather*} F_n(x_2;w_1)=\tau(k)F_n\big(k^{-1}x_2;k^*w_1\big)\rho(k)^{-1}. \end{gather*}
Therefore for $k={\rm e}^{-t\hbar}$ we have
\begin{gather}\label{invariance_Fn}
F_n(x_2;w_1)=\chi_\rho\big({\rm e}^{t\hbar}\big)\chi_\tau\big({\rm e}^{-t\hbar}\big) F_n\big({\rm e}^tx_2;\overline{{\rm e}^{-t}}w_1\big),
\end{gather}
where the notations are the same as in the previous subsection.
Now we fix $t>0$, and let $|x_2|_\infty<{\rm e}^{-t}<1$. Then for $v\in V$ we have
\begin{gather*}
 \left|\left(F_n\left(x_2;\overline{\frac{\partial}{\partial x_1}}\right)f(x_1),v\right)_\tau\right| =\left|\left(\chi_\rho\big({\rm e}^{t\hbar}\big)\chi_\tau\big({\rm e}^{-t\hbar}\big)
F_n\left({\rm e}^tx_2;\overline{{\rm e}^{-t}\frac{\partial}{\partial x_1}}\right)f(x_1),v\right)_\tau\right| \\
\qquad{} =\frac{\chi_\rho({\rm e}^{t\hbar})\chi_\tau({\rm e}^{-t\hbar})}{n!}\left|\left\langle {\rm e}^{-nt}
\left(y_1\,\middle|\overline{\frac{\partial}{\partial x_1}}\right)_{\fp^+_1}^nf(x_1),
\hat{\rK}\big({\rm e}^tx_2;y_1\big)^*v\right\rangle_{\hat{\rho},y_1}\right| \\
\qquad{} \le \chi_\rho\big({\rm e}^{t\hbar}\big)\chi_\tau\big({\rm e}^{-t\hbar}\big)\frac{{\rm e}^{-nt}}{n!}
\left\Vert\left(y_1\,\middle|\overline{\frac{\partial}{\partial x_1}}\right)_{\fp^+_1}^nf(x_1)\right\Vert_{\hat{\rho},y_1}
\big\Vert\hat{\rK}\big({\rm e}^tx_2;y_1\big)^*v\big\Vert_{\hat{\rho},y_1}.
\end{gather*}
Next we estimate
$\ds \left\Vert\left(y_1\,\middle|\overline{\frac{\partial}{\partial x_1}}\right)_{\fp^+_1}^nf(x_1)\right\Vert_{\hat{\rho},y_1}$.
For $x_1,y_1\in D_1$, we take $R>0$ such that $x_1+R{\rm e}^{\sqrt{-1}\theta}y_1\in D_1$ for any $\theta\in\BR$. Then
\begin{align*}
\left|\left(y_1\,\middle|\,\overline{\frac{\partial}{\partial x_1}}\right)_{\fp^+_1}^nf(x_1)\right|_\rho
&=\left|\left.\frac{{\rm d}^n}{{\rm d}t^n}\right|_{t=0}f(x_1+ty_1)\right|_\rho
=\left|\frac{n!}{2\pi\sqrt{-1}}\oint_{|z|=R}\frac{f(x_1+zy_1)}{z^{n+1}}{\rm d}z\right|_\rho \\
&=\left|\frac{n!}{2\pi\sqrt{-1}}\int_0^{2\pi}\frac{f\big(x_1+R{\rm e}^{\sqrt{-1}\theta}y_1\big)}{R^{n+1}{\rm e}^{\sqrt{-1}(n+1)\theta}}
\sqrt{-1}R{\rm e}^{\sqrt{-1}\theta}{\rm d}\theta\right|_\rho \\
&=\left|\frac{n!}{2\pi}\int_0^{2\pi}\frac{f\big(x_1+R{\rm e}^{\sqrt{-1}\theta}y_1\big)}{R^n{\rm e}^{\sqrt{-1}n\theta}}{\rm d}\theta\right|_\rho\\
& \le \frac{n!}{R^n}\max_{\theta\in[0,2\pi]}\big|f\big(x_1+R{\rm e}^{\sqrt{-1}\theta}y_1\big)\big|_\rho.
\end{align*}
Now fix $0<s<1$, let $|x_1|_\infty<s$, and we set $\ds R:=\frac{s-|x_1|_\infty}{|y_1|_\infty}$. Then we have
\begin{gather*} \big|x_1+R{\rm e}^{\sqrt{-1}\theta}y_1\big|_\infty\le |x_1|_\infty+R|y_1|_\infty=s, \end{gather*}
and therefore
\begin{align*}
\left|\left(y_1\,\middle|\,\overline{\frac{\partial}{\partial x_1}}\right)_{\fp^+_1}^nf(x_1)\right|_\rho
&\le n!\left(\frac{|y_1|_\infty}{s-|x_1|_\infty}\right)^n\sup_{|\xi_1|_\infty\le s}|f(\xi_1)|_\rho \\
&\le n!\left(\frac{1}{s-|x_1|_\infty}\right)^n\sup_{|\xi_1|_\infty\le s}|f(\xi_1)|_\rho
\end{align*}
holds. Hence we have
\begin{gather}
 \left\Vert\left(y_1\,\middle|\overline{\frac{\partial}{\partial x_1}}\right)_{\fp^+_1}^nf(x_1)\right\Vert_{\hat{\rho},y_1}^2 \notag \\
\qquad{} =C_\rho\int_{D_1} \left(\rho\big(B(x_1)^{-1}\big)\left(y_1\,\middle|\,\overline{\frac{\partial}{\partial x_1}}\right)_{\fp^+_1}^nf(x_1)
,\left(y_1\,\middle|\,\overline{\frac{\partial}{\partial x_1}}\right)_{\fp^+_1}^nf(x_1)\right)_\rho h_1(y_1)^{-p_1}{\rm d}y_1 \notag \\
\qquad{} \le C_\rho\int_{D_1} \big|\rho\big(B(x_1)^{-1}\big)\big|_{\rho,\mathrm{op}}h_1(y_1)^{-p_1}{\rm d}y_1
\sup_{y_1\in D_1}\left|\left(y_1\,\middle|\,\overline{\frac{\partial}{\partial x_1}}\right)_{\fp^+_1}^nf(x_1)\right|_\rho^2 \notag \\
\qquad{} \le C_\rho^{\prime 2}(n!)^2\left(\frac{1}{s-|x_1|_\infty}\right)^{2n}\sup_{|\xi_1|_\infty\le s}|f(\xi_1)|_\rho^2. \label{higherdiff}
\end{gather}
Therefore we have
\begin{gather*}
 \left|\left(F_n\left(x_2;\overline{\frac{\partial}{\partial x_1}}\right)f(x_1),v\right)_\tau\right| \\
 \le C_\rho'\chi_\rho\big({\rm e}^{t\hbar}\big)\chi_\tau\big({\rm e}^{-t\hbar}\big)
\left(\frac{{\rm e}^{-t}}{s-|x_1|_\infty}\right)^n\sup_{|\xi_1|_\infty\le s}|f(\xi_1)|_\rho
\big(\big\langle\hat{\rK}\big({\rm e}^tx_2;\cdot\big)^*,
\hat{\rK}\big({\rm e}^tx_2;\cdot\big)^*\big\rangle_{\hat{\rho}}v,v\big)_\tau^{1/2}.
\end{gather*}
As a consequence, if we assume ${\rm e}^{-t}<s-|x_1|_\infty$, we get
\begin{gather*}
 \left|\left((\cF_{\tau\rho}f)(x),v\right)_\tau\right|\le\sum_{n=0}^\infty
\left|\left(F_n\left(x_2;\overline{\frac{\partial}{\partial x_1}}\right)f(x_1),v\right)_\tau\right| \\
 \le C_\rho'\chi_\rho\big({\rm e}^{t\hbar}\big)\chi_\tau\big({\rm e}^{-t\hbar}\big)
\sum_{n=0}^\infty \left(\frac{{\rm e}^{-t}}{s-|x_1|_\infty}\right)^n\!\sup_{|\xi_1|_\infty\le s}\!|f(\xi_1)|_\rho
\big(\big\langle\hat{\rK}\big({\rm e}^tx_2;\cdot\big)^* ,
\hat{\rK}\big({\rm e}^tx_2;\cdot\big)^*\big\rangle_{\hat{\rho}}v,v\big)_\tau^{1/2} \\
 =C_\rho'\chi_\rho\big({\rm e}^{t\hbar}\big)\chi_\tau\big({\rm e}^{-t\hbar}\big)\frac{s-|x_1|_\infty}{s-|x_1|_\infty-{\rm e}^{-t}}\sup_{|\xi_1|_\infty\le s}|f(\xi_1)|_\rho
\big(\big\langle\hat{\rK}\big({\rm e}^tx_2;\cdot\big)^*,
\hat{\rK}\big({\rm e}^tx_2;\cdot\big)^*\big\rangle_{\hat{\rho}}v,v\big)_\tau^{1/2}.
\end{gather*}
This estimate holds if $0<{\rm e}^{-t}<s<1$, $|x_1|_\infty<s-{\rm e}^{-t}$ and $|x_2|_\infty<{\rm e}^{-t}$, and thus
\begin{gather*}
 ((\cF_{\tau\rho}f)(x),v )_\tau=\sum_{n=0}^\infty\left(F_n\left(x_2;\overline{\frac{\partial}{\partial x_1}}\right)f(x_1),v\right)_\tau
\end{gather*}
converges absolutely and uniformly on every compact subset in $\{|x_1|_\infty<s-{\rm e}^{-t},|x_2|_\infty<{\rm e}^{-t}\}$ for any $0<{\rm e}^{-t}<s<1$,
and hence also does on every compact subset in $\{|x_1|_\infty+|x_2|_\infty<1\}$.
By Theorem \ref{conti_ext}, $\cF_{\tau\rho}$ is a continuous operator from $\cO_\rho(D_1,W)$ to $\cO_\tau(D,V)$,
and therefore $\cF_{\tau\rho} f(x)$ must extend to whole $D$.
\end{proof}
